\documentclass[11pt]{article}
\usepackage[margin=1in]{geometry}
\usepackage[T1]{fontenc}
\usepackage{lmodern,microtype}
\usepackage{amsmath,amssymb,amsthm,mathtools}
\usepackage{booktabs,tabularx,longtable,array,enumitem}
\usepackage{xcolor,listings,fancyhdr}
\usepackage[colorlinks=true,linkcolor=blue!45!black,citecolor=blue!45!black,urlcolor=blue!45!black]{hyperref}
\usepackage{bookmark}
\hypersetup{pdftitle={iPRA: coverage before cancellation and smaller proof kernels},pdfauthor={Research prepared for Vladimir Reshetnikov with ChatGPT},pdfsubject={Explicit iPRA presentations, interpretation proofs, and reproducible checker measurements}}
\setlength{\emergencystretch}{3em}
\setlength{\parskip}{0.3em}
\setlist{itemsep=3pt,topsep=5pt}
\newtheorem{theorem}{Theorem}[section]
\newtheorem{lemma}[theorem]{Lemma}
\newtheorem{proposition}[theorem]{Proposition}
\newtheorem{corollary}[theorem]{Corollary}
\theoremstyle{definition}\newtheorem{definition}[theorem]{Definition}
\newcommand{\ipra}{\mathrm{iPRA}}
\newcommand{\Q}{\mathsf Q}
\newcommand{\D}{\mathsf D}
\newcommand{\file}[1]{\expandafter\nolinkurl\expandafter{\detokenize{#1}}}
\lstset{language=Python,basicstyle=\ttfamily\scriptsize,columns=fullflexible,
 keepspaces=true,breaklines=true,breakatwhitespace=false,showstringspaces=false,
 frame=single,rulecolor=\color{black!20},backgroundcolor=\color{black!2},
 xleftmargin=3pt,xrightmargin=3pt,aboveskip=8pt,belowskip=8pt}
\pagestyle{fancy}\fancyhf{}
\fancyhead[L]{\small iPRA: coverage and proof kernels}
\fancyhead[R]{\small 10 November 2026}
\fancyfoot[C]{\thepage}
\setlength{\headheight}{14pt}

\title{\textbf{iPRA: coverage before cancellation,\\
standard-numeral rigidity, and smaller proof kernels}}
\author{Research prepared for Vladimir Reshetnikov\\with ChatGPT}
\date{10 November 2026}

\begin{document}
\maketitle
\begin{abstract}
This investigation produces a new coverage theorem for the scalar
self-seeded arithmetic in the supplied drafts. Without assuming successor
injectivity, the weak theory proves that every element is zero or a successor,
that successor has no fixed points, and that every fixed standard positive
numeral has exactly its intended predecessor. With injectivity retained,
the predecessor construction needs only two exceptional inputs and no joint
coverage/predecessor induction. A separate logical theorem identifies one
native paired binder over an equation with the explicitly prenex
existential-universal quantifier-free class, explaining precisely why
weakening the induction guard is substantive.

Complete proof checkers for the inherited initialized-iterator presentation
are also reduced: the recommended explicit-list program has 539 normalized
Python bytes; a compact equivalent has 531. A distinct generic profile has
289 AST nodes. The current comparison envelope is 542 bytes, 261 tokens,
and 291 AST nodes, attained by different programs. No new program beats
all three coordinates, and the token minimum is unchanged. The arithmetic
lemma is not assigned the sizes of these inherited-arithmetic checkers.

\textbf{Unresolved target.} This is not a proof of an absolutely simplest
axiomatization or the requested significant complete-checker breakthrough.
Successor injectivity in the weak self-seeded theory remains unproved here;
neither its full iPRA bi-interpretability nor a countermodel is claimed.
The archive contains complete sources, attributed equivalence proofs,
executable certificates, measurements, and negative controls.
\end{abstract}
\newpage
\setcounter{tocdepth}{1}
\tableofcontents
\newpage

\section{Result, comparison, and scope}
\label{sec:result}
The audit freezes the 20 direct ZIP archives in
\file{VladimirReshetnikov/ProveIt}, under \file{docs/drafts}, at this
commit:\cite{source}
\begin{center}\small
\texttt{b9c4d74a0018bdafc5bc06e755a0145f4ed2b8d6}
\end{center}
The static audit checks each archive's Git blob identity and length,
recurses through 35 distinct ZIP contents, and measures 484 distinct
Python source byte strings. There are 2,744 source occurrences under the
deduplicated-archive traversal. At or below any previously reported
candidate threshold, 48 distinct source files require classification:
10 complete theorem checkers and 38 support programs. The latter are
build scripts, test runners, translators, measurers, or integrity tools.
They cannot be counted as standalone theorem checkers.

A final read-only check of commit \texttt{89f5bdfe} found that the
archives had been curated into
\file{Logic/PrimitiveRecursiveArithmetic/ConciseIPRA/Research}.
All 2,769 retained research files match original archive members by Git
blob identity, including all 1,605 Python files and all 187 TeX files.
Only the project README and attributes file are new. All three
comparison-frontier programs remain present and unchanged; the migration
does not alter this comparison.\cite{maintained}
The exact hashes and remeasurements are recorded in
\file{evidence/latest_snapshot_recheck.json}.

The observed coordinate minima in both inspected snapshots are
\[
  (\text{bytes},\text{tokens},\text{AST nodes})=(542,261,291).
\]
The first two occur in different neutral-major profiles; the third occurs
in the earlier lexical-loop profile.\cite{neutral,selfseed}
They are observed minima within this fixed corpus, hence upper bounds on
an unknown optimum. They are not mathematical lower bounds on all possible
theories or implementations.

\begin{center}\small
\begin{tabular}{lrrrr}
\toprule
Whole standalone module & Bytes & Tokens & AST & Gzip-9\\
\midrule
Current byte leader &542&263&309&343\\
Current token leader &578&261&295&346\\
Current AST leader &642&263&291&353\\
\midrule
New explicit-list checker &539&266&308&340\\
New compact-list checker &531&266&310&335\\
New generic-loop checker &613&263&289&349\\
\bottomrule
\end{tabular}
\end{center}
The byte improvement is $11/542\approx2.03\%$ for the compact profile.
The AST improvement is $2/291\approx0.69\%$ for a different program.
The explicit-list program is recommended for review and ordinary use:
its parsing recursion follows finite JSON trees. The generic profile
has a malformed-string divergence qualification under an ideal unlimited
stack, explained in Section~\ref{sec:kernel}. These programs implement
the same initialized-iterator arithmetic. Their size reductions compose
existing neutral-major, lexical-scope, and control-flow techniques;
they are not presented as new arithmetic axioms.

There are three logically separate contributions:
\begin{enumerate}
\item The new theorem about the weaker self-seeded arithmetic in
Section~\ref{sec:coverage-no-cancellation}; it does not complete an
injectivity-free iPRA foundation.
\item The binder-depth expressivity calculation in
Section~\ref{sec:depth-one-boundary}; it is a precise induction-guard
diagnostic, not a universal checker-size lower bound.
\item Smaller complete presentations of the inherited, fully
iPRA-equivalent calculus in Sections~\ref{sec:axioms} and
\ref{sec:kernel}; these satisfy full-domain bi-interpretability through
the included inherited bridge.
\end{enumerate}

\section{The complete calculus implemented by the measured checkers}
\label{sec:axioms}
This section explicitly fixes every primitive numerical schema and
inference pattern. The construction is the inherited Choice-DI calculus
in formula-context form, with a new composition of certificate techniques.
The complete mathematical theory must not be confused with the weaker
self-seeded theory investigated in the next section.

\subsection{Numerical and formula syntax}
There is one numerical sort. Terms are
\[
 t::=x\mid0\mid St\mid I(n,a;u.h).
\]
Only $u$ in $h$ is bound by the iterator. Every update is fixed finite
earlier-stage term syntax; all other variables are passive numerical
parameters. Equivalently, introduce one first-order symbol $I_h$ for
each such finite description. There are no numerical function codes
and no uniformly total evaluator.

The two computation schemas are
\begin{align}
 I(0,a;u.h)&=a, & I(Sn,a;u.h)&=h[I(n,a;u.h)/u].
 \tag{I0, IS}\label{eq:choice-computations}
\end{align}
Native formulas are
\[
 A::=(s=t)\mid\Q_y A\mid\D_{x,y}A.
\]
The designated absurdity is $0=S0$. The intended ordinary readings are
\[
 (\Q_yB)^\circ=B^\circ(0)\to\forall yB^\circ(y),\qquad
 (\D_{x,y}B)^\circ=
 \exists x\bigl(B^\circ(x,0)\to\forall yB^\circ(x,y)\bigr).
\]
These are interpretation clauses, not additional primitive connectives.
The $Q$ binder is conservative shorthand for unused-witness $D$;
its explicit presence makes the checker interface smaller.

\subsection{All primitive inference patterns}
Contexts consist of reusable native formulas. Substitutions avoid capture.
The complete rule inventory, in addition to I0 and IS, is:
\begin{enumerate}[label=\textbf{\arabic*.}]
\item \textbf{Hypothesis.} Infer $\Gamma\vdash A$ when $A\in\Gamma$.
\item \textbf{Transport.} From $\Gamma\vdash s=t$ and
$\Gamma\vdash H(s)$, infer $\Gamma\vdash H(t)$, for any fixed native
formula template $H$.
\item \textbf{Closure elimination.} From $\Gamma\vdash\Q_yB(y)$
and $\Gamma\vdash B(0)$, infer $\Gamma\vdash B(t)$.
\item \textbf{Paired elimination.} From $\Gamma\vdash\D_{x,y}B(x,y)$
and $\Gamma,\Q_yB(a,y)\vdash A$, infer $\Gamma\vdash A$.
The parameter $a$ is fresh for $\Gamma,A$ and the passive matrix data.
The displayed local assumption is discharged.
\item \textbf{Case introduction.} From
$\Gamma,B(t,0)\vdash B(t,Sz)$, infer
$\Gamma\vdash\D_{x,y}B(x,y)$.
\item \textbf{Equation-progress introduction.} From
$\Gamma,E(t,z)\vdash E(t,Sz)$, infer
$\Gamma\vdash\D_{x,y}E(x,y)$, only when $E$ is a literal numerical
equation. An equation below a binder is not a literal equation.
\end{enumerate}
In the last two rules $z$ is fresh for the retained context, witness,
and fixed matrix parameters; the displayed assumption is discharged.
Weakening, exchange, contraction, renaming, cut, and numerical
substitution are admissible by induction on finite derivations.
No generalization rule allowing a retained parameter to change during
an induction is used.

No further numerical or logical axioms are primitive here. Reflexivity,
successor injectivity, numerical coverage, the ordinary intuitionistic
connectives, and full quantifier-free induction are derived. The
appendices supply the construction in a noncircular order.

\subsection{Why this is fully equivalent to iPRA}
The equation-progress rule, an unused witness in $D$, and its elimination
derive contextual literal-equation induction. Case introduction similarly
gives numerical cases. Computation and transport give reflexivity:
if $r=I(0,t;u.0)$, use $r=t$ as both transport premises with the
motive $v=t$ to obtain $t=t$.

The derived selector $C(v,a,b)=I(v,a;u.b)$ permits formula gating and
selection. Appendix~\ref{sec:choice-introduction} explicitly reconstructs
intuitionistic implication, conjunction, disjunction, both quantifiers,
and falsity, including both formula round trips. The scalar predecessor
is constructed using only equation induction and numerical cases in
Appendix~\ref{sec:hybrid-bootstrap}. Equality reflection then yields
quantifier-free induction; pairing compiles every fixed general PR
description. Internal recurrence uniqueness identifies both numerical
return translations. The isomorphism of the original domain with its
double interpretation has graph $x=y$ in both theories.

Thus the theory implemented by each selected checker is equiconsistent
and bi-interpretable with full first-order iPRA. This statement concerns
all native formulas and retained contexts, including quantified ones,
and the entire numerical domain. It is stronger than equality of standard
numerical functions or equiconsistency alone. The underlying mathematical
bridge is inherited and explicitly attributed; the new checker composition
does not change it.

\section{Coverage before cancellation in the scalar self-seeded basis}
\label{sec:coverage-no-cancellation}

This section isolates a mathematical simplification of the supplied
scalar self-seeded presentation. It is not a complete new basis below
the measured checker frontier. In particular, it does not justify
deleting successor injectivity from that presentation.

Use ordinary intuitionistic first-order logic with equality, terms
$0,S,\epsilon$, and a staged scalar recursor $F(n;v.h)$ with equations
\[
 F(0;v.h)=S(h[0/v]),\qquad
 F(Sn;v.h)=h[F(n;v.h)/v].
\]
The remaining nonlogical assumptions, besides contextual induction on
literal equations, are explicitly
\[
 0=S0\to\bot,\qquad
 \epsilon(a,b)=0\lor\epsilon(a,b)=1,\qquad
 \epsilon(a,b)=1\leftrightarrow a=b.
\]
All free numerical variables in these schemas are universally closed. Here $1=S0$, and each recursor body is an earlier-stage
source term. Neither successor injectivity nor numerical coverage is
assumed in the result below.

Put $D(c,a)=F(c;v.a)$ and $Z(c)=D(c,0)$. Thus
\[
 D(0,a)=Sa,\quad D(Sn,a)=a,\quad Z(0)=1,\quad Z(Sn)=0.
\]
If $Sn=0$, these equations give $0=1$, so every successor is nonzero.
For Boolean $b,c$, the source terms
\[
 \neg_Bb=\epsilon(b,0),\qquad
 b\wedge_Bc=\epsilon(D(b,\epsilon(c,0)),0)
\]
have their usual truth tables; define $\vee_B$ by De Morgan. Only
$0\ne1$ and successor nonzeroness are used in these finite cases.
The usual external syntax recursion assigns a Boolean characteristic
to each fixed quantifier-free formula. Literal-equation induction
therefore gives contextual quantifier-free induction, without
successor cancellation.

Define the following fixed source terms:
\[
 \begin{split}
 B_a(v)&=\epsilon(v,0)\vee_B\epsilon(Sv,a),\\
 H_a(v)&=D(B_a(v),v),\\
 U(j,a)&=F(j;v.H_a(v)),\qquad p_*(a)=U(a,a).
 \end{split}
\]
The update holds the state at zero or immediately before $a$, and
increments it in the remaining case. Since $B_a(0)=1$,
$H_a(0)=0$ and $U(0,a)=1$. With $a$ fixed, quantifier-free induction gives
\begin{equation}
 U(j,a)=Sj\quad\lor\quad S(U(j,a))=a.
 \label{eq:protected-orbit}
\end{equation}
At zero the first alternative holds. A state satisfying the second
alternative is fixed by $H_a$. Otherwise the first alternative advances
from $Sj$ to $S(Sj)$: in that branch the current state is a successor
and hence nonzero. No separate induction asserting nonzeroness of the
whole orbit is needed.

After this induction has been proved with its parameter fixed,
specialize $j=m$ and $a=S^2m$. Both alternatives of
\eqref{eq:protected-orbit} imply $S(U(m,S^2m))=S^2m$. The state is
then fixed. Exactly two further computation steps give
\begin{equation}
                     S(p_*(S^2m))=S^2m.
 \label{eq:protected-right-inverse}
\end{equation}
This is a right-inverse equation only; no cancellation of $S$ has
taken place.

\paragraph{Whole-domain coverage theorem.}
The theory with the assumptions just listed proves
\[
                     n=0\ \lor\ \exists m\,n=Sm.
\]
Indeed, apply quantifier-free induction to
\[
 C(n):\quad n=0\ \lor\ n=1\ \lor\ S(p_*(n))=n.
\]
The zero case is immediate. In the successor step, $n=0$ gives
$Sn=1$. If $n=1$, apply \eqref{eq:protected-right-inverse} with
$m=0$. In the third case $n=S(p_*(n))$, whence
$Sn=S^2(p_*(n))$; apply \eqref{eq:protected-right-inverse} with
$m=p_*(n)$ and rewrite. This proves $C(n)$ for every $n$.
Its second and third alternatives provide the witnesses $0$ and
$p_*(n)$ respectively. The proof covers nonstandard elements and
never assumes that an arbitrary element is reached by a standard
finite successor iteration.

\subsection{No fixed points and rigidity of each standard numeral}

The same weak theory proves more than coverage, although the following
results still fall short of arbitrary successor cancellation.

\paragraph{Successor has no fixed point.}
Let $h(v)=\epsilon(v,0)$. This source update has no fixed point.
Indeed, $h(v)=v$ and the Boolean range of $h$ give $v=0$ or $v=1$;
the two cases contradict $h(0)=1$ and $h(1)=0$, respectively.
If $Sa=a$, the recursor equation would give
\[
 F(a;v.h)=F(Sa;v.h)=h(F(a;v.h)),
\]
contrary to that fact. Thus $Sa\ne a$ for every $a$, without using
successor injectivity or coverage. In particular, $1\ne2$.

\paragraph{The only predecessor of one is zero.}
Put
\[
 h(v)=\neg_B\epsilon(v,2),\qquad A(n)=F(n;v.h).
\]
The distinctness of $0,1,2$ gives $h(0)=1$, $A(0)=2$, and
$A(1)=h(2)=0$. The image of $h$ is contained in $\{0,1\}$, and
\[
                         h(v)=0\quad\longleftrightarrow\quad v=2.
\]
Quantifier-free induction proves $A(n)=2\leftrightarrow n=0$.
The zero case holds, whereas at every successor both sides are false:
$A(Sn)=h(A(n))$ is zero or one, and $Sn\ne0$.
If $Sn=1$, then $h(A(n))=A(Sn)=A(1)=0$; hence $A(n)=2$ and $n=0$.
The converse is immediate. Therefore
\[
                              Sn=1\quad\longleftrightarrow\quad n=0.
\]

\paragraph{Finite-prefix rigidity metatheorem.}
For every fixed standard integer $k\ge1$, there is a finite derivation
in the weak theory of
\begin{equation}
 Sx=\overline{k}\quad\longleftrightarrow\quad x=\overline{k-1}.
 \label{eq:standard-predecessor}
\end{equation}
Here bars denote fixed numerals. This is an external metatheorem,
not a derivation of a single formula quantifying over $k$.

The proof is an external induction on $k$. Carry as induction data
the distinctness of the numerals $0,\ldots,k+1$ and
\eqref{eq:standard-predecessor} for all smaller positive indices.
The initial distinctness of $0,1,2$ follows above. Define the following
finite numeral-valued source update:
\[
 h_k(v)=
 \begin{cases}
 \overline{i-1},&v=\overline{i},\quad 1\le i\le k-1,\\
 \overline{k-1},&v=\overline{k+1},\\
 \overline{k},&\text{otherwise}.
 \end{cases}
\]
This display abbreviates an actual source term, not an additional
recursion schema over a graph-defined function. To construct it,
for each fixed $r=1,\ldots,k$ form the Boolean characteristic $b_r(v)$
of the finite condition that the displayed value is at least $r$.
Explicitly, $b_r$ is the Boolean negation of the disjunction of
$\epsilon(v,\overline{i})$ for $1\le i\le\min(r,k-1)$, together
with $\epsilon(v,\overline{k+1})$ when $r=k$. Empty disjunctions
mean zero; all index comparisons here are metalevel comparisons
between fixed integers, not assumed arithmetic order predicates.
Start with $t_0(v)=0$ and set
\[
                  t_r(v)=D(\neg_B b_r(v),t_{r-1}(v)).
\]
Then $h_k=t_k$: each true condition contributes one successor, and
each false condition leaves the value unchanged. The construction
uses exactly $k$ external syntax steps and the already available
finite Boolean gates. It is not asserted to provide a single source
term with variable $k$.

Write $A_k(n)=F(n;v.h_k)$ and introduce the distinct marker numerals
\[
 q_0=\overline{k+1},\qquad q_j=\overline{k-j}\quad(1\le j\le k).
\]
Since $h_k(0)=\overline{k}$, the initial state is $q_0$. Direct finite
computation gives $A_k(\overline{j})=q_j$ for $0\le j\le k$.
The displayed update has no output $q_0$, and for $1\le j\le k$ it
satisfies
\[
                       h_k(v)=q_j\quad\longleftrightarrow\quad v=q_{j-1}.
\]
Quantifier-free induction therefore gives
$A_k(n)=q_0\leftrightarrow n=0$. For $1\le j<k$, an external
induction on $j$, using numerical coverage and the already proved
predecessor uniqueness for $\overline{j}$, gives
\[
                     A_k(n)=q_j\quad\longleftrightarrow\quad n=\overline{j}.
\]
For clarity, in the successor case $n=Sm$ the left side is equivalent
to $A_k(m)=q_{j-1}$, hence to $m=\overline{j-1}$. The right side has
exactly the same condition by the previous predecessor theorem.
The zero case has two false sides.

Now $Sx=\overline{k}$ implies
$h_k(A_k(x))=A_k(\overline{k})=q_k$, so $A_k(x)=q_{k-1}$ and
$x=\overline{k-1}$. This proves the next predecessor theorem.
To extend numeral distinctness to $k+2$, a putative equality
$\overline{k+2}=\overline{j}$ with $1\le j\le k$ reduces, by the
newly available predecessor theorems, to
$\overline{k+1}=\overline{j-1}$, already impossible. Equality with
zero is excluded by successor nonzeroness, and equality with $k+1$
by the no-fixed-point lemma. This closes the external induction.

Consequently no element other than the intended previous numeral can
map to any fixed standard positive numeral. This does not rule out
collisions over nonstandard targets: replacing the fixed construction
above by a uniform variable-index argument would require an additional
theorem that has not been supplied.

\paragraph{A simpler predecessor when cancellation is retained.}
If successor injectivity is retained, as in the supplied complete
self-seeded theory, \eqref{eq:protected-right-inverse} gives
$p_*(S^2m)=Sm$. Define the total single-valued finite-case graph
\[
 P(n)=
 \begin{cases}
 0,&n=0\text{ or }n=1,\\
 p_*(n),&n\ne0\text{ and }n\ne1.
 \end{cases}
\]
Its graph uses only quantifier-free source guards and source value
terms. Finite common refinements therefore eliminate each fixed
quantifier-free formula containing $P$, preserving the available
induction rule. The zero equation is immediate. For $P(Sn)=n$,
use the already derived coverage to split $n=0$ or $n=Sm$.
The first gives $P1=0$. In the second, $S^2m$ is neither zero nor
one, and the displayed equation gives $P(S^2m)=Sm=n$.

Thus only the values at zero and one need correction, rather than
the three exceptional values in the supplied stopped-counter
construction. Coverage is proved first, and the subsequent
predecessor proof no longer needs a joint coverage/predecessor
induction. The remaining open issue is precise: can the weaker
theory prove $p_*(S^2m)=Sm$, or otherwise prove $Sx=Sy\to x=y$?
The proof above establishes neither statement.

\section{The exact logical strength visible at one native binder}
\label{sec:depth-one-boundary}

This section returns to the complete initialized-$I$ native calculus of
Section~\ref{sec:axioms}; its selector, predecessor, and PR compiler have
been established. Those resources have not been claimed for the weaker
self-seeded $F$ theory of the preceding section.

This observation locates a concrete boundary that a proposed simplification
of the induction guard must respect. It is a theorem about the established
arithmetic interpretation of the native syntax, not a non-bi-interpretability
theorem for every conceivable unrelated interpretation.

Write $C(0,a,b)=a$ and $C(Sn,a,b)=b$ for the derived numerical selector,
and $P(0)=0$, $P(Sn)=n$ for the derived predecessor. All formulas below may
have arbitrary fixed numerical parameters. A formula is called
$\Sigma_2^{\mathrm{qf}}$ here if it has the form
$\exists x\,\forall y\,\theta(x,y,\vec p)$ with quantifier-free $\theta$.
Several variables in a block may be packaged by the already derived
primitive-recursive pairing operations. This notation specifies an actual
prenex class; it does not appeal to classical prenex normalization of an
arbitrary intuitionistic formula.

\begin{proposition}[One-binder expressivity]
Under the whole-domain interpretation in iPRA, the class of native formulas
$D_{x,z}E(x,z,\vec p)$ with $E$ a literal equation is exactly the class
$\Sigma_2^{\mathrm{qf}}$, up to provable equivalence. Similarly, native
$Q_zE(z,\vec p)$ formulas are exactly the corresponding
$\Pi_1^{\mathrm{qf}}$ class, up to provable equivalence.
\end{proposition}

\begin{proof}
The forward direction follows from
\[
 D_{x,z}E\equiv
 \exists x\bigl(E(x,0)\to\forall z\,E(x,z)\bigr)
 \leftrightarrow
 \exists x\,\forall z\bigl(E(x,0)\to E(x,z)\bigr).
\]
The final matrix is quantifier-free. The displayed interchange of
implication and universal quantification is intuitionistically valid,
because $z$ is absent from $E(x,0)$. Dropping the outer existential gives
the analogous statement for $Q$.

Conversely, choose the iPRA characteristic term $r(x,y,\vec p)$ for a
given quantifier-free matrix, with
$r(x,y,\vec p)=0\leftrightarrow\theta(x,y,\vec p)$. Expand $C$, $P$, and $r$
into their fixed compiled initialized-iterator source terms, so the following
formula is literally an equation of the native language. Define the literal
equation
\[
 E(x,z,\vec p):\quad
 C\bigl(z,0,r(x,Pz,\vec p)\bigr)=0.
\]
Its zero instance is provable. Its successor instance is equivalent to
$r(x,y,\vec p)=0$, since $P(Sy)=y$. Numerical coverage therefore proves
\[
 \forall z\,E(x,z,\vec p)
 \leftrightarrow \forall y\,\theta(x,y,\vec p).
\]
Since $E(x,0,\vec p)$ is provable, its occurrence as the antecedent in
the native $D$ reading disappears by intuitionistic reasoning. Thus
$D_{x,z}E$ has precisely the required existential-universal meaning.
The same construction with no existential coordinate establishes the
$Q$ claim.
\end{proof}

\begin{corollary}[A one-binder progress extension already adds a full
prenex induction class]
Suppose the native equation-progress introduction rule is extended to
matrices $B(n,\vec p)$ that consist of one $D$ binder over a literal
equation. Over the existing native theory, this extension derives every
induction instance for $\Sigma_2^{\mathrm{qf}}$ formulas. Extending the
rule to one $Q$ binder over an equation derives the analogous
$\Pi_1^{\mathrm{qf}}$ induction instances.
\end{corollary}

\begin{proof}
Let $B(n,\vec p)$ be the native translation supplied by the proposition.
Suppose $\Gamma\vdash B(0,\vec p)$ and
$\Gamma,B(z,\vec p)\vdash B(Sz,\vec p)$, with the usual fresh $z$.
Apply the extended progress rule using an unused existential coordinate.
It yields $D_{u,n}B(n,\vec p)$. Paired elimination exposes
$Q_nB(n,\vec p)$, and closure elimination with the given base instance
proves $B(t,\vec p)$ for an arbitrary term $t$. The internal equivalence
from the proposition transports this derivation to the original prenex
formula. The argument for $Q$ matrices is identical.
\end{proof}

Consequently, retaining a test that merely says ``at most one binder''
does not preserve the stated literal-equation induction restriction.
Such a change supplies induction for this additional syntactic class;
equivalence with the original rule would require a separate elimination
or conservativity theorem. No claim here identifies every prenex
induction extension with full HA, and no strength separation for
arbitrary alternative bi-interpretations is inferred from this calculation.

\subsection{A uniform lower bound on what larger depths can express}

The one-binder equivalence extends to a useful inclusion at all depths.
The qualification is important: no intuitionistic prenex-normal-form
theorem is being silently used for nested native implications.
Let $\delta(E)=0$ for equations and
$\delta(Q_yA)=\delta(D_{x,y}A)=1+\delta(A)$.

For a fresh numerical parameter $z$, define $G_zA$ by replacing every
equation $s=t$ in the fixed finite formula $A$ by
\[
                 C(z,0,s)=C(z,0,t),
\]
and retaining its native binders. Capture is avoided by renaming the
bound variables first. This transformation does not change $\delta$.
It has the internal properties
\[
             \vdash G_0A,\qquad G_{Su}A\leftrightarrow A.
\tag{Gate}\label{eq:depth-gate}
\]
For equations these are computation identities. For a $Q$ or $D$ root,
the first assertion follows inductively because its entire matrix is
true, using the zero witness for $D$; the second follows from the same
interpreted binder applied to equivalent matrices. This reasoning uses
the already recovered full logical interface, rather than assuming a
new matrix-congruence rule during an earlier arithmetic bootstrap.

\begin{theorem}[Alternating-prefix encoding at depth $k$]
For every $k\geq0$, every explicitly prenex formula
\[
 \exists x_1\forall y_1\cdots\exists x_k\forall y_k\,
 \theta(\vec x,\vec y,\vec p),
\]
with quantifier-free $\theta$, has an iPRA-provably equivalent native
formula of depth at most $k$. The translation is uniform by external
recursion on the displayed finite prefix. In particular, allowing
equation-progress introduction on all matrices of depth at most $k$
derives induction for this entire prenex class.
\end{theorem}

\begin{proof}
For $k=0$, replace the quantifier-free formula by its primitive-recursive
characteristic equation. For the induction step suppose $F(x,y,\vec p)$
already translates the remaining prefix at depth at most $k-1$. Put
\[
             \mathcal B(F)
       :=D_{x,z}\,G_zF(x,Pz,\vec p).
\]
By \eqref{eq:depth-gate}, its zero matrix is true and its successor
matrix at $z=Sy$ is equivalent to $F(x,y,\vec p)$. Numerical coverage
then gives
\[
             \mathcal B(F)
       \leftrightarrow\exists x\,\forall y\,F(x,y,\vec p).
\]
Its depth is at most $k$, establishing the recursive construction.
For induction, repeat the derivation in the preceding corollary with
this translated matrix. The equivalences commute with the passive
induction parameter and its zero and successor substitutions.
\end{proof}

Blocks of several numerical variables can be packaged into one tuple
using the available pairing and projections. Dummy blocks can be added
because the numerical domain is inhabited. Consequently the inclusion
covers the usual explicitly prenex $\Sigma_{2k}$ class, and for $k\ge1$ the
$\Pi_{2k-1}$ class by an unused first existential block. It does not
claim that every intuitionistic formula of a conventional arithmetical
hierarchy has been prenexed without additional assumptions.

At depths greater than one, only this encoding inclusion is claimed
intuitionistically. The native reading contains implications whose
antecedents may themselves be quantified; moving such quantifiers
through implication is not a general intuitionistic equivalence.
Under classical logic the usual prenex manipulations give the familiar
$\Sigma_{2k}$ upper bound for a $D$ root of depth $k$, but that classical
observation must not be substituted for an iPRA proof. The new inclusion
and the exact depth-one calculation already suffice to audit the
literal-equation guard.



\section{The complete proof kernels and their invariants}
\label{sec:kernel}
The recommended complete program is printed here. It imports only the
standard JSON parser. The theorem producer, tests, arithmetic metatheory,
and corpus audit are not trusted imports.
\begin{lstlisting}
import json

def c(a, g=[[]], k=0):
    i = lambda a, t=None, d=0: [i(b, t, d + 5 - len(a)) for b in a] if isinstance(a, list) else a if k <= a < d else d % 2 and i(t, d=3)
    if i(a) in g:
        return
    m, s, u, t = q()
    c(m, g, k)
    a = m in ([s, u], not s and [t] * 3) and i(t, u) == a and i(t, s) or a
    h = m * 3 or (u and [[s and [s], u, t], i(t, [s, u, t], 3) if s else u]) or (a := i(*a, t, 2))
    k -= 1
    c(i(a, [k]), g + [i(h, (m or len(a) % 2 == u) and k)], k)
q = json.load(open(0)).pop
c(q())
print('ACCEPT')
\end{lstlisting}
The more compact list test in \file{code/checker_bytes.py} replaces
\texttt{isinstance(a, list)} by \texttt{a * 0 == []}. For JSON values,
only arrays satisfy that test; unsupported multiplication or a failed
numerical comparison rejects the input. The explicit test is retained
in the recommended program for readability. A fully expanded standalone
version is supplied as \file{code/readable_checker.py}.

\subsection{Public contract and sorted expression encoding}
Input is one finite JSON array whose last element is the goal. Earlier
four-field packets $[m,s,u,t]$ are consumed in reverse order. Acceptance
requires normal process exit and exactly \texttt{ACCEPT} followed by a
newline. Exceptions, nonzero exits, resource failures, and other output
are rejection. Unused certificate prefixes are permitted. Completeness
is the usual adequate-resource claim about finite canonical proofs;
there is no assertion that fixed hardware checks arbitrarily large proofs.

The inherited representation is sorted: the numerical and formula
interpretations of an array differ.
\begin{center}\small
\begin{tabularx}{\linewidth}{@{}cXX@{}}
\toprule
Arity & Numerical interpretation & Formula interpretation\\\midrule
0 & Zero & The theorem $0=0$\\
1 & Successor & Paired binder $D$\\
2 & Zero; annotations ignored & Equality\\
3 & Scalar iterator $I$ & Closure $Q$, last child its matrix\\
$\ge4$ & Zero; annotations ignored & The theorem $0=0$\\
\bottomrule
\end{tabularx}
\end{center}
Canonical forms include $[n]$ for successor, $[n,a,h]$ for iteration,
$[s,t]$ for an equation, $[B]$ for $D$, and $[B,B,B]$ for $Q$.
Formula root depth is zero and numerical root depth is three. Under
$p$ protected numerical slots the depths are $2p$ and $3+2p$.
Coordinates $0,1,2$ denote numerical zero; coordinates $3+2j,4+2j$
name protected slot $j$. Admitted negative labels name rigid parameters.

The common child-depth increment is $5-r$ at an array of arity $r$.
A successor child is read beneath two dummy zero slots. The counter
and initializer of an iterator are beneath one dummy zero slot; its
update is beneath the actual state slot. The same increment therefore
implements the different binding conventions. A scalar escaping a
protected numerical environment is the substitution hole, whose image
is checked at numerical root depth three. Formula scalar padding has
the inherited truth meaning. This is a representation of finite terms,
not a new arithmetic operation or a total evaluator for numerical codes.

\subsection{Lexical scopes and the crucial allocation order}
At a call $c(a,g,k)$, the opener closes over that call's fresh-name floor
$k$. A major subproof receives the same context and floor value. Its
private allocation cannot change the caller's floor. Context extension
constructs a new list in the recursive programs. In the loop profile,
the context is an immutable tuple and \texttt{+=} rebinds it. The default
list in the recursive programs is never mutated.

The current goal is checked at the old floor. Raw hypothesis construction
and witness opening occur before decrementing $k$; every recursive-image
or witness check triggered there uses that old floor. The final hypothesis
opening may admit fresh names, discharged by the zero-instance argument
below. Different completed major
branches may reuse names because their assumptions cannot escape. A
finite raw certificate may use fresh fractional numerical labels; these
are treated as a finite family of fresh parameters, as in the inherited
raw-input proof. Canonical certificates use only integer coordinates.

The order is essential. The delivered four-packet attack in
\file{tests/future_witness_capture.json} targets
\[
 D_{x,y}\bigl[I(y,0;v.x)=I(y,0;v.y)\bigr].
\]
Its zero equation is true. At every successor $y$, its equation says
$x=y$. Thus its ordinary meaning requires a single number equal to
every positive number, which is false already at $1$ and $2$.
An unsafe mutation that allocates the new name before opening the
witness accepts this formula: the witness can then be chosen as the
successor of the induction eigenparameter. The corrected programs
reject it. This is a concrete reason to retain the construction-before-
allocation order, not merely a formatting preference.

\subsection{Neutral major premises and folded branches}
The initial hypothesis \texttt{[]} denotes the already derived theorem
$0=0$. Every non-hypothesis command checks its major $m$ unconditionally.
A canonical command without a substantive major supplies \texttt{[]};
its proof closes immediately and allocates no name. The inherited
list-valued-context invariant excludes numeric scalar majors from a
finite successful checking trace. In particular, scalar zero cannot
reach the matcher as a false-valued major. This is why the potentially
dangerous Python \texttt{False} alternative in the matcher is harmless.
The complete raw argument is reproduced in the final appendix.

For a surviving nonempty list $m$, the list $m*3$ is nonempty. For a
computation command, the constructed hypothesis is a two-element list,
also nonempty. Consequently the original three-way choice of major
export, computation, and introduction is equivalently implemented by
\[
 \texttt{m*3 or (u and COMPUTATION) or INTRODUCTION}.
\]
No semantic branch is removed. The AST profile folds only the outer
major choice. Its loop replaces a tail recursive continuation with a
context-preserving iteration; it does not change an inference rule.

In an introduction, both $m$ and $u$ are false-valued. The mode comparison
is with $u$: zero selects progress for a literal equation, while an
empty array selects the zero case. No false-valued JSON mode equals the
odd arity of a meaningful compound matrix. Replacing $u$ by the
unrestricted source field $s$ would permit compound induction; the
archive contains a targeted forbidden-induction certificate for that error.

These changes are compositions of documented existing techniques. Their
particular complete source files were absent from the frozen corpus,
but no foundational novelty is inferred from that fact.

\subsection{Computation supply and missing literal-zero cases}
The direct computation interface uses truthy syntactic initializers,
and a truthy predecessor for a successor instance. This is not a
numerical nonzeroness test. The finite term
\[
                     Z_*:=I(S(S0),S0;v.0)
\]
has a directly permitted proof of $Z_*=0$. For an excluded computation,
replace its offending zero predecessor and initializer positions
simultaneously by a fresh placeholder in a fixed equation $E(w)$.
Then $E(Z_*)$ is directly permitted; transport along $Z_*=0$ gives
$E(0)$. One transport and its two computation premises repair the
leaf. Thus unrestricted I0 and IS are recovered without an arithmetic
oracle or an extra axiom.

A generated raw hypothesis may contain newly admitted names and need
not be true uniformly in those names. The sound rule is
\[
 \vdash H(\vec0),\qquad\Gamma,H(\vec b)\vdash A
 \quad\Longrightarrow\quad\Gamma\vdash A,
 \qquad\vec b\notin\operatorname{FV}(\Gamma,A).
\]
Substitute zero for all the finitely many new names in the continuation,
rename later eigenparameters first, and cut in the proved zero instance.
The old-floor validation is exactly what licenses this discharge.

\subsection{Soundness and canonical completeness}
Induct on a finite successful trace. Hypothesis closure uses a retained
formula. A recursive major is proved in the caller's old context.
Endpoint matching realizes transport or closure elimination. A unary
paired major exports its closure at a fresh witness; other surviving
arities export only truth padding. Introduction retains exactly the
permitted case or literal-equation progress hypothesis. Computation
supply uses the preceding zero-instance discharge argument. Finally
cut away the initial derived truth. This yields a derivation in the
native calculus of Section~\ref{sec:axioms}.

Conversely, start with a finite native derivation, normalize the
excluded computation leaves as above, and alpha-rename its parameters
with a branch-local name supply. Major subproofs start at the caller's
floor, and continuations allocate their own next name. Canonical absent
majors are empty arrays. Folding the nonempty branches chooses the same
hypothesis, and looping the tail continuation has the same effect.
The resulting finite certificate is accepted. The decoded canonical
theorems are therefore exactly those of the complete native calculus;
its full-domain iPRA interpretations transfer unchanged.

This argument does not claim identical accepted languages for all
redundant malformed JSON spellings. In particular, the generic opener
admits some iterable truth aliases, while folded multiplication may
reject them before a continuation. That only removes redundant raw
encodings and leaves canonical completeness intact.

\subsection{The generic profile's qualification}
The generic AST profile attempts iterable traversal and catches failure.
A nonempty malformed string can cause repeated traversal of its own
one-character strings. Real CPython ultimately rejects through its
stack limit; an ideal unlimited-stack interpreter can diverge. Thus
289 AST nodes is not claimed as the size of a total abstract parser on
all malformed JSON. The recommended explicit-list program and the
compact-list variant recurse structurally through finite trees and
reject unsupported leaf types. The arithmetic and metatheoretic
claims use finite successful traces and finite canonical proofs.

The written proofs have not been formalized in Lean or Rocq. The
regression suite tests the implementation and concrete counterexamples;
it does not prove the universal soundness or bi-interpretability theorem.


\section{Evidence, reproducibility, and limitations}
\label{sec:evidence}
All recorded Python measurements use CPython 3.12.14 and zlib 1.3.2.
The measurer removes genuine docstrings, applies \texttt{ast.unparse},
and appends a newline. It then counts UTF-8 bytes, significant Python
tokens, and every AST node of the whole standalone module, including
JSON input and the verdict driver. The common Python/JSON runtime is
excluded, matching the inspected drafts' convention. There is no hidden
theory-specific import, embedded binary theorem engine, or decompression
trick in a selected checker.

Gzip is reported with \texttt{compresslevel=9} and \texttt{mtime=0},
including its wrapper. It is an informational metric. Raw and normalized
gzip lengths are both retained. For the three normalized compact sources
these lengths coincide. The readable expansion has 2,307 raw bytes,
1,463 normalized bytes, 295 tokens, and 348 AST nodes; its raw gzip size
is 961 bytes and its normalized gzip size is 546 bytes.

The checker suite has 49 local certificate cases and 46 global-reference
cases. Together with malformed-input cases it checks 756 complete-process
expectations across the four new programs and three frozen baselines,
under ordinary and optimized Python. Six additional expectations run
three deliberately invalid implementation variants in both modes. All pass with zero
mismatches. The mutations exercise compound-progress selection, context
leakage from a subproof, and future-variable capture by premature name
allocation. Their targeted certificates are rejected by the
corrected programs and accepted by the corresponding mutants.

The independent arithmetic diagnostics evaluate the displayed source
Boolean operations and protected orbits. They check 361 orbit/parameter
pairs, 19 right-inverse inputs, 21 coverage/predecessor inputs, 21
fixed-point/marker inputs, six fixed-prefix indices, and 63 finite-update
source cases. Seven targeted bad controls detect missing zero protection,
missing stopping, omitted finite exceptions, and corrupted marker
operations. A separate branching-successor example illustrates the
invalidity of cancelling $S$ from a right-inverse identity alone. It is
explicitly \emph{not} a model of all the weak theory's recursor and
induction schemas.

Two independent assistant workers reviewed the coverage and fixed-prefix
arguments, with the detailed logical review retained in the archive.
This is internal independent checking, not external human peer review
or proof-assistant certification. The tests do not establish a theorem
about nonstandard elements; that conclusion rests on the written
fixed-parameter induction arguments.

\subsection{Reproduction commands and files}
From the extracted archive directory, run
\begin{verbatim}
python3 verify_manifest.py
python3 reproduce.py
python3 build.py
\end{verbatim}
The first verifies distributed file identities. The second remeasures the
complete modules and frozen corpus, regenerates the certificates, checks
the process expectations, and runs the arithmetic diagnostics. The third
builds the standalone TeX source and PDF with shell escape disabled.
The tests require only Python's standard library. Building the report
requires a standard pdfLaTeX installation.

\file{evidence/snapshot.json} fixes the source commit and archive blob
identities. \file{evidence/corpus_metrics.json} records all 484 distinct
source measurements and locations; \file{evidence/boundary_classification.json}
records the 48 relevant role classifications. Exact frozen Python
sources are supplied under \file{evidence/source_corpus/}, so the static
remeasurement needs no network. To repeat the full archive-identity audit
against a local checkout of the pinned repository, run
\begin{verbatim}
python3 audit_corpus.py /path/to/ProveIt/docs/drafts
\end{verbatim}
No inspected corpus program is executed by that audit. Only the selected
and explicitly attributed test programs are executed during reproduction.

The complete compact and readable checkers are under \file{code/};
the four-packet witness-capture attack and regenerated named proofs are
under \file{tests/}. \file{reference/} retains the measured incumbents,
the original inherited mathematical sources, and the upstream MIT-0
license. Attribution and content hashes are recorded in
\file{evidence/mathematical_provenance.json} and the source manifest.
The adapted arithmetic appendix makes its scalar prerequisites explicit
before general PR compilation, avoiding a possible circular reading of
the original exposition.

The complete deliverable supports the specific new lemmas, small source
reductions, and inherited full-iPRA interpretation stated here. It does
not establish a single all-coordinate record, an absolute minimum,
or an injectivity-free complete self-seeded axiomatization.


\section{Further research questions}
\begin{enumerate}
\item Can the weak self-seeded theory prove successor injectivity?
The proved identity is $S(p_*(S^2m))=S^2m$. A derivation of
$p_*(S^2m)=Sm$ would close a central gap; simply cancelling the outer
$S$ assumes exactly what is missing.
\item Can the fixed-standard-index marker construction be internalized
uniformly without assuming a general primitive-recursive compiler?
Every fixed numeral target has a unique predecessor, but passing from
that external family to one variable-index theorem is not automatic.
\item If cancellation is not derivable, can one construct a model of
\emph{all} staged recursor equations and equation-induction instances
with a successor collision? A branching successor graph alone does not
meet those obligations.
\item Can a single complete, structurally terminating checker beat the
current 261-token minimum and the new 531-byte and 289-node results
simultaneously? Merely mixing coordinates of different programs does
not answer this question.
\item Can a new certificate invariant remove a whole rule or validation
operation while preserving the entire consequence relation? The concrete
scope and compound-induction attacks in this archive are useful regression
targets for such a proposal.
\item Can the new coverage, standard-target, and binder-depth proofs be
formalized in Lean or Rocq using only the exact source schemas? Numerical
testing and an informal stronger induction principle would not constitute
that formalization.
\end{enumerate}

\appendix
\section{The ordinary intuitionistic logic used in the arithmetic arguments}
\label{sec:ordinary-rules}
The native calculus has exactly the primitive rules in
Section~\ref{sec:axioms}; it derives the following interface. The weaker
self-seeded theory in Section~\ref{sec:coverage-no-cancellation} instead
uses this ordinary interface explicitly. It is therefore charged as part
of that theory, not treated as free hidden machinery.

Formulas are built from numerical equations and $\bot$ with
$\wedge,\vee,\to,\forall,\exists$. The rules are single-conclusion
intuitionistic natural deduction:
\[
\frac{A\quad B}{A\wedge B}\quad
\frac{A\wedge B}{A}\quad\frac{A\wedge B}{B}\qquad
\frac{A}{A\vee B}\quad\frac{B}{A\vee B}\quad
\frac{A\vee B\quad[A]\ \vdots\ C\quad[B]\ \vdots\ C}{C},
\]
\[
\frac{[A]\ \vdots\ B}{A\to B}\qquad
\frac{A\to B\quad A}{B}\qquad
\frac{\bot}{A},
\]
\[
\frac{A(a)}{\forall xA(x)}\quad
\frac{\forall xA(x)}{A(t)}\qquad
\frac{A(t)}{\exists xA(x)}\quad
\frac{\exists xA(x)\quad[A(a)]\ \vdots\ B}{B}.
\]
In universal introduction, $a$ is fresh for the undischarged context.
In existential elimination, $a$ is fresh for that context, the conclusion
$B$, and the passive parameters of the quantified formula. Brackets
mark the hypotheses discharged by the inference. There is the hypothesis
rule and the usual reusable-context structural behavior. Equality has
reflexivity $t=t$ and capture-free transport from $s=t,A(s)$ to $A(t)$
for every fixed formula template. These imply the usual symmetry,
transitivity, and function congruence laws. Negation abbreviates
$A\to\bot$; biconditional abbreviates the conjunction of both implications.
Neither excluded middle for arbitrary formulas nor double-negation
elimination is a rule.

For reference, the contextual literal-equation induction schema is
\[
\frac{\Gamma\vdash s(0,\vec p)=t(0,\vec p)\qquad
 \Gamma,s(z,\vec p)=t(z,\vec p)\vdash s(Sz,\vec p)=t(Sz,\vec p)}
 {\Gamma\vdash s(a,\vec p)=t(a,\vec p)}.
\]
The parameter $z$ is fresh for the retained context, target, and fixed
parameters. This restriction is part of the schema. The derived
quantifier-free induction used in the new coverage proof is obtained
by translating each fixed quantifier-free formula into its Boolean
characteristic equation, rather than by silently allowing a larger
primitive induction rule.

\section{Reconstructing the intuitionistic logical interface}
\label{sec:choice-introduction}
This appendix reproduces the relevant part of the supplied Choice-DI development, with its Q-context form used in the main text. It is included to make the inherited logical interpretation explicit.\cite{source}

\subsection{Induction and numerical cases separate cleanly}

\begin{lemma}[Contextual equation induction]
\label{lem:choice-ei}
The progress choice, DE, and Use derive
\[
\frac{\Gamma\vdash E(0)\qquad
      \Gamma,E(z)\vdash E(Sz)}
     {\Gamma\vdash E(t)}
\tag{EI}\label{eq:choice-ei}
\]
for every equation $E$, with the usual freshness condition.
\end{lemma}
\begin{proof}
Choose an unused witness coordinate $u$. Progress DI introduces
$\mathsf D_{u,y}E(y)$ from the given step. DE opens the clause
$\langle E(0);y.E(y)\rangle$. Supply its antecedent using the given
base proof and use the clause at $t$. Close DE.
\end{proof}

\begin{lemma}[Cases, equality, and binder congruence]
\label{lem:choice-preinduction}
The case choice derives, for arbitrary native formulas,
\[
\frac{\Gamma\vdash F(0)\qquad\Gamma\vdash F(Sz)}
     {\Gamma\vdash F(t)}
\tag{NC}\label{eq:choice-nc}
\]
with the usual freshness condition. Equality reflexivity,
symmetry, transitivity, numerical congruence, and uniform
matrix interderivability under $\mathsf D$ are derivable
without using the progress choice.
\end{lemma}
\begin{proof}
For cases, use case DI with an unused witness coordinate and
matrix $F(y)$, ignoring its retained base in the successor
premise. DE opens $\langle F(0);y.F(y)\rangle$; the given zero
proof and Use at $t$ finish.

For reflexivity put $r=I(0,t;u.0)$. Use the computation equation
$r=t$ as both transport premises with motive $v=t$ to obtain
$t=t$. The usual explicit equality motives give symmetry,
transitivity, and numerical congruence.

For binder congruence, open the old paired formula at a fresh
witness $a$ and introduce the new one by case DI at $a$.
Convert its retained base to the old base, use the opened
old clause at the successor, and convert back. Ordinary cut
closes the construction, and the reverse conversion gives
the converse. No induction is used.
\end{proof}

The primitive progress choice is still restricted to equations.
Allowing it at arbitrary compound matrices would, by the
argument for \eqref{eq:choice-ei}, supply induction for
arbitrary formulas and change the intended theory.

\subsection{The entire ordinary logical interface}
\label{subsec:choice-logic}

Here is a predecessor-free reconstruction of the ordinary logical
operations. It makes explicit that the case choice supplies the
logical interface independently of progress induction.

Let $C(v,a,b):=I(v,a;u.b)$ with $u$ unused. Computation gives
$C(0,a,b)=a$ and $C(Sv,a,b)=b$. Define the gate on fixed finite
formula syntax by
\[
 G_v(s=t):=(C(v,0,s)=C(v,0,t)),\qquad
 G_v(\mathsf D_{x,y}B):=\mathsf D_{x,y}G_vB,
\]
renaming bound variables away from the control. Equality,
binder congruence, and external structural induction give
\[
 \vdash G_0A,\qquad G_{Sv}A\dashv\vdash A,\qquad
 A\dashv\vdash\mathsf D_{u,y}G_yA
 \quad(u,y\notin\operatorname{FV}(A)).
\tag{Gate and padding}\label{eq:choice-gate}
\]
For padding, case DI proves the forward direction. DE, the
zero-gate theorem, and Use at $1$ prove the reverse direction.

Define a selector between two arbitrary formulas by
\[
 M_v(s=t,u=w):=(C(v,s,u)=C(v,t,w))
\]
and, with aligned fresh paired variables, by
\[
\begin{aligned}
 M_v(\mathsf D_{x,y}A,\mathsf D_{x,y}B)
   &:=\mathsf D_{x,y}M_v(A,B),\\
 M_v(E,\mathsf D_{x,y}B)
   &:=\mathsf D_{x,y}M_v(G_yE,B),\\
 M_v(\mathsf D_{x,y}A,E)
   &:=\mathsf D_{x,y}M_v(A,G_yE),
\end{aligned}
\]
where $E$ is an equation. Every call decreases maximum input
binder depth. Padding and congruence prove
$M_0(A,B)\dashv\vdash A$ and
$M_{Sv}(A,B)\dashv\vdash B$.

Put $\delta(v)=C(v,1,0)$ and $p(z)=I(z,0;u.\delta(u))$.
Only the fixed computations $p(1)=1$ and $p(2)=0$ are needed.
The full logical macro inventory is
\[
\begin{aligned}
 A\Rightarrow_*B&:=\mathsf D_{u,z}M_z(A,B),\\
 A\wedge_*B&:=\mathsf D_{u,z}G_zM_{p(z)}(A,B),\\
 \exists_*x\,A(x)&:=\mathsf D_{x,y}G_yA(x),\\
 A\vee_*B&:=\exists_*x\,M_x(A,B),\\
 \forall_*x\,A(x)&:=
        A(0)\wedge_*\mathsf D_{u,z}G_zA(z),\\
 \bot_*&:=(0=1).
\end{aligned}
\tag{Logical macros}\label{eq:choice-macros}
\]
Newly bound variables are fresh; $u$ is unused.

Implication introduction obtains its assumption from the zero
selector and supplies the consequent at successors, then uses
case DI. Elimination opens its clause and uses it at $1$.
Existential introduction uses its witness and the successor
gate; elimination opens at a fresh witness, supplies the
zero-gate theorem, and uses the clause at $1$.

Disjunction introduction uses witnesses $0$ and $1$.
For elimination, discharge the two branch proofs as
$A\Rightarrow_*F$ and $B\Rightarrow_*F$. NC applied to
$M_z(A,B)\Rightarrow_*F$ gives this implication at every
control: the two endpoint cases follow from the branch
implications. Existential elimination provides $M_a(A,B)$
and the implication gives $F$.

Given $A$ and $B$, selector endpoints and NC give $M_v(A,B)$
at an arbitrary control. Substitute $v=p(Sz)$ and use the
successor gate for conjunction introduction. For projections,
open its zero-gated clause and specialize at $1$ and $2$;
the two fixed computations give $B$ and $A$, respectively.

These rules prove existential numerical coverage:
\[
     (x=0)\vee_*\exists_*v\,(x=Sv).
\tag{Coverage}\label{eq:choice-coverage}
\]
Use NC. Zero takes the left injection of reflexivity, and a
successor takes the right injection with the evident witness.
Coverage gives equality cases even when the retained context
contains $x$; its elimination adds $x=0$ or $x=Sv$ rather
than silently changing that retained context.

Universal introduction follows from a fresh-parameter proof
by substituting zero and successors and forming the two
conjuncts. For elimination, take its $A(0)$ conjunct and
open the paired conjunct. Its base is a zero-gate theorem.
Coverage of the target chooses the zero instance from
$A(0)$ or a successor instance from the opened clause
and gate removal.

Finally, under $0=1$, computation and transport give
$s=C(0,s,t)=C(1,s,t)=t$ for arbitrary numerical terms.
Every equation follows, and external induction on formula
syntax with case DI gives every native formula. This
supplies falsity elimination. No excluded middle for
quantified formulas has entered the construction.

\subsection{Whole-domain realizations and both formula composites}

For a standard first-order reading, let $T_{\rm eq}$ have
intuitionistic equality logic, the same numerical signature
and computation equations, numerical coverage, separation
$(0=1)\to\bot$, and contextual equation induction.
Its interpretation of the native syntax is
\[
 (s=t)^\circ=(s=t),\qquad
 (\mathsf D_{x,y}B)^\circ
   =\exists x\bigl(B^\circ(x,0)\to\forall y B^\circ(x,y)\bigr),
\]
\[
 \langle A;y.B\rangle^\circ
       =A^\circ\to\forall y B^\circ.
\]
The domain is the entire numerical domain and interpreted
equality is actual equality.

Case DI is sound in this realization by numerical coverage.
Progress DI is sound by contextual equation induction under
the conditional base. Use, DE, transport, and computation
have their displayed ordinary meanings. Conversely, the
native macros derive every ordinary logical rule,
\eqref{eq:choice-ei} derives equation induction, and
\eqref{eq:choice-coverage} derives coverage. Separation
translates to the identity implication
$(0=1)\Rightarrow_*(0=1)$.

For the essential paired formula composite,
\[
 \mathsf D_{x,y}B
 \dashv\vdash
 \exists_*x\bigl(B(x,0)\Rightarrow_*\forall_*y\,B(x,y)\bigr).
\tag{Paired composite}\label{eq:choice-paired-composite}
\]
Forwards, DE opens its clause at a fresh $a$; under $B(a,0)$
that clause supplies $B(a,y)$ at a fresh arbitrary $y$.
Introduce the universal, implication, and existential
macros and close DE. Backwards, eliminate the existential
at a fresh $a$. Introduce the paired formula with witness
$a$ by case DI: the retained base, supplied implication,
and universal elimination at $Sz$ give its successor premise.

The ordinary interpretations of the six logical macros
have their ordinary round trips. For conjunction, its true
base and the instances $1,2$ recover both operands;
the converse uses selector cases. The universal has all
successor instances in its paired conjunct and zero in
its first conjunct; coverage gives exactly the ordinary
universal. Implication, disjunction, and existential
quantification use their selector or gate endpoints.
Falsity uses the explicit ordinary separation bridge.
External induction on finite formulas, together with
\eqref{eq:choice-paired-composite}, gives both formula
composites. Clauses correspond in both directions to
$A\Rightarrow_*\forall_*yB$ by the derived logical rules.

The numerical composites are literal identities with
internally provable graph $x=y$. No standard-numeral
restriction, retract, quotient, or equality reflection
assumption is being inserted at this stage. The arithmetic bridge in the following appendices identifies $T_{\rm eq}$ with
full iPRA; the arithmetic appendices below establish that bridge for the complete native rule inventory stated in the main text.


% Adapted with attribution from the pinned unparenthesized one-hole archive.
\section{Internal scalar predecessor bootstrap}
\label{sec:hybrid-bootstrap}

This scalar bootstrap is inherited from the supplied scalar development~\cite{erasure} and is reproduced to keep the scalar checker self-contained. It is not a new contribution of this paper. The proof uses a factor map fixing odd states and a saturated-iteration construction of predecessor.

The available arithmetic terms are $0$, $S$, and
\[
 I(n,a;u.h),\qquad
 I(0,a;u.h)=a,\qquad
 I(Sn,a;u.h)=h[I(n,a;u.h)/u].
\]
The update binds only $u$.  All other variables in it are fixed numerical
parameters.  We use equality reasoning, the already reconstructed
zero/successor cases, and induction on a single literal equation.
Every induction below keeps its other parameters fixed.  Predecessor,
successor injectivity, pairing, and induction on arbitrary
quantifier-free formulas are not assumptions of the argument.

\subsection{Elementary terms and two finite cycles}

Write $1=S0$ and use larger numerals as successor abbreviations.  Define
\[
 C(n,a,b)=I(n,a;u.b),\qquad
 Z(n)=C(n,1,0),\qquad
 a+b=I(b,a;u.Su).
\]
The update variable of $C$ is absent from its body.  Consequently
\[
 C(0,a,b)=a,\qquad C(Sn,a,b)=b.
\]
Numerical cases show that $Z(n)$ is $0$ or $1$.
The usual commutative-monoid equations for addition follow by literal
equation induction.  For example, $0+b=b$ and $Sa+b=S(a+b)$ are
proved by induction on $b$; these yield commutativity by a further
equation induction, and associativity follows by induction on the last
summand.  We use these equations to reorder finite sums.  The notation
$2a$ means $a+a$ and $3a$ means $a+a+a$; no multiplication operation
is needed for this bootstrap.

Put
\[
 N(n)=I(n,0;u.Z(u)),\qquad
 j(u)=Z(u)+2N(u),\qquad
 M(n)=I(n,0;u.j(u)).
\]
The range of $N$ is contained in $\{0,1\}$: its zero value is $0$,
and its successor value is a value of $Z$.  This uses numerical cases
on its counter, not induction on a disjunction.  Hence
\begin{equation}
 N(n)+Z(N(n))=1,\qquad
 Z(Z(N(n)))=N(n),\qquad N(n+2)=N(n).
 \label{eq:hybrid-parity}
\end{equation}
Here the last equation follows from two successor computations and
the preceding equation.  Literal-equation inductions give
\begin{equation}
 N(2n)=0,\qquad N(2n+1)=1,\qquad N(3n)=N(n).
 \label{eq:hybrid-parity-scaling}
\end{equation}
For the first equation the successor step adds two to the argument of
$N$; the second follows by one computation.  In the third, increasing
$n$ adds three to the left argument and therefore applies $Z$ once to
both sides of the induction equation.

Finite computations now give
\[
 j(0)=1,\qquad j(1)=2,\qquad j(2)=0.
\]
The proof below needs just these three values and
$M(Sn)=j(M(n))$; it does not assume a division algorithm or a theorem
characterizing $M$ as remainder modulo three.

\subsection{A ramp with a global scaling identity}

Define
\[
 G(n)=I(n,0;u.u+M(u)+1).
\]
The central invariant is
\begin{equation}
 M(G(n))=N(n).
 \label{eq:hybrid-ramp-invariant}
\end{equation}
Its zero case is immediate.  For its successor step, retain the
single induction equation and split the Boolean value $N(n)$.
If it is $0$, the next $G$ state is $G(n)+1$; its $M$ value is
$j(0)=1$, equal to $N(Sn)$.  If it is $1$, the next $G$ state
is $G(n)+2$; two computations of $M$ change its value from $1$
to $2$ and then to $0$, again equal to $N(Sn)$.  This proves
the invariant by literal-equation induction.

It follows that
\[
 G(Sn)=G(n)+N(n)+1,\qquad G(n+2)=G(n)+3.
\]
Equation induction therefore yields
\begin{equation}
 G(2n)=3n.
 \label{eq:hybrid-ramp-double}
\end{equation}
A separate literal-equation induction proves
\begin{equation}
 G(3n)=3G(n)+N(n).
 \label{eq:hybrid-ramp-triple}
\end{equation}
For completeness, its successor calculation is
\begin{align*}
 G(3Sn)
 &=G(3n)+N(3n)+4
   =3G(n)+2N(n)+4,\\
 3G(Sn)+N(Sn)
 &=3G(n)+3N(n)+3+Z(N(n)).
\end{align*}
The expressions agree in both Boolean cases for $N(n)$.
The base is $0=0$.  In particular,
\eqref{eq:hybrid-ramp-triple} was proved internally; it was not
inferred from an intended floor-function interpretation of $G$.

Set
\[
 H(u)=C(N(u),G(u),u).
\]
The crucial choice is to retain $u$ when its parity is $1$.
The preceding equations give
\begin{equation}
 H(2u)=3u,\qquad H(1)=1,\qquad H(3u)=3H(u).
 \label{eq:hybrid-factor-map}
\end{equation}
To check the last identity, use $N(3u)=N(u)$ and split the Boolean
$N(u)$.  At $0$, both sides are $3G(u)$ by
\eqref{eq:hybrid-ramp-triple}; at $1$, both are $3u$.
This argument does not require that an arbitrary number has already
been represented as $2v$ or $2v+1$.

\subsection{A comparator from scalar iteration}

For every fixed update $h$, equation induction on $n$, keeping $m,a$
and all parameters fixed, proves the semigroup identity
\[
 I(m+n,a;h)=I(n,I(m,a;h);h).
\]
The base is the zero computation and the step applies $h$ to the
induction equation. This identity will also be used in the later
pairing construction.


For a fixed update $h$, equation induction on $y$, with its initializer
$a$ and all parameters of $h$ fixed, proves the rotation identity
\begin{equation}
 I(Sy,a;h)=I(y,h(a);h).
 \label{eq:hybrid-rotation}
\end{equation}
Both sides have base $h(a)$, and their successor equations apply the
same $h$ to the previous equality.  Here and below $I(y,a;h)$ is
an abbreviation for one fixed finite scalar term, not quantification
over functions.

By \eqref{eq:hybrid-factor-map}, another equation induction gives
\begin{equation}
 I(y,3a;H)=3I(y,a;H).
 \label{eq:hybrid-iterator-scaling}
\end{equation}
Its base is $3a=3a$.  In the step the fixed-parameter induction
equation rewrites the new left side as
\[
 H(3I(y,a;H))=3H(I(y,a;H)).
\]
The initializer parameter stays $a$ throughout this induction.
Equation induction also gives $I(y,1;H)=1$, since $H(1)=1$.

Define
\[
 E(x)=I(x,1;u.2u),\qquad
 L(x,y)=N\bigl(I(y,E(x);H)\bigr).
\]
Then the following entire recursion specification is provable:
\begin{equation}
 L(0,y)=1,\qquad L(Sx,0)=0,\qquad L(Sx,Sy)=L(x,y).
 \label{eq:hybrid-comparator}
\end{equation}
The first equation uses $E(0)=1$ and the fixed-point identity for
$H$.  The second uses $E(Sx)=2E(x)$ and $N(2a)=0$.
For the third,
\begin{align*}
 L(Sx,Sy)
 &=N\bigl(I(Sy,2E(x);H)\bigr)\\
 &=N\bigl(I(y,H(2E(x));H)\bigr)\\
 &=N\bigl(I(y,3E(x);H)\bigr)\\
 &=N\bigl(3I(y,E(x);H)\bigr)\\
 &=N\bigl(I(y,E(x);H)\bigr)=L(x,y).
\end{align*}
The successive steps use rotation, $H(2a)=3a$, iterator scaling,
and $N(3a)=N(a)$.  Thus $L$ has the comparison equations at all
arguments internally.  Its standard-model meaning as the indicator
of $x\le y$ is motivation, not a premise of the following argument.

\subsection{Saturated iteration gives predecessor}

We now need only the three equations
\eqref{eq:hybrid-comparator}.  Define
\[
 B_m(u)=C(L(Su,m),u,Su),\qquad
 W(t,m)=I(t,0;u.B_m(u)).
\]
The subscript in $B_m$ records a fixed parameter; $B_m$ is an open
term abbreviation, not a function variable.

Numerical cases on $e$ prove the selector equation
\[
 C(e,Sa,Sb)=S C(e,a,b).
\]
Equivalently one may use equation induction on $e$, whose induction
hypothesis is unused.  Therefore
\begin{align}
 B_{Sm}(Su)
 &=C(L(SSu,Sm),Su,SSu)\notag\\
 &=C(L(Su,m),Su,SSu)=S B_m(u),
 \label{eq:hybrid-conjugacy}\\
 B_{Sm}(0)&=C(L(S0,Sm),0,S0)=C(L(0,m),0,S0)=S0.
 \label{eq:hybrid-conjugacy-base}
\end{align}

Literal-equation induction on $t$, with $m$ fixed, proves
\begin{equation}
 W(St,Sm)=S W(t,m).
 \label{eq:hybrid-saturated-shift}
\end{equation}
The base is \eqref{eq:hybrid-conjugacy-base}.  The successor step is
\begin{align*}
 W(SSt,Sm)
 &=B_{Sm}(W(St,Sm))\\
 &=B_{Sm}(S W(t,m))\\
 &=S B_m(W(t,m))=S W(St,m).
\end{align*}
Only the second line uses the current induction hypothesis, at the
same fixed $m$.

Two subsequent equation inductions give
\begin{equation}
 L(Sn,n)=0,\qquad W(n,n)=n.
 \label{eq:hybrid-diagonals}
\end{equation}
The first uses the comparator recurrence.  The second uses
\eqref{eq:hybrid-saturated-shift} in its successor step.  The
substitution $t=m=n$ is made into that already proved universally
parameterized equation; no parameter of an active induction hypothesis
is changed.  Consequently
\begin{equation}
 W(Sn,n)=B_n(n)=C(L(Sn,n),n,Sn)=n.
 \label{eq:hybrid-extra-step}
\end{equation}

Finally define the finite scalar term
\begin{equation}
 P(m)=I\bigl(m,0;u.C(L(SSu,m),u,Su)\bigr).
 \label{eq:hybrid-predecessor-term}
\end{equation}
Its zero equation is immediate.  To justify its successor equation
without assuming a separate extensionality rule for iterator bodies,
put
\[
 V(t,m)=I\bigl(t,0;u.C(L(SSu,Sm),u,Su)\bigr).
\]
The comparator recurrence identifies this update pointwise with
$B_m(u)$.  Equation induction on $t$, with $m$ fixed, gives
$V(t,m)=W(t,m)$: their base values are $0$, and the successor step
uses the pointwise update equation and then the induction equation.
Now \eqref{eq:hybrid-extra-step} yields
\[
 \boxed{\qquad P(0)=0,\qquad P(Sm)=V(Sm,m)=W(Sm,m)=m.\qquad}
\]

\subsection{Dependency boundary}

All terms in this section are built from $0,S,I$ by finite substitution.
All inductions have a single equation as their invariant.  Numerical
cases and equality transport are the already reconstructed logical
resources.  Successor injectivity may now be obtained by applying $P$
to an assumed equation $Sx=Sy$.  Equality tests, reflection of
quantifier-free formulas, pairing, and the recovery of general
primitive recursion can then be developed in the established order.
They were not used in order to manufacture predecessor.

This argument is a written internal derivation in the stated theory.
The finite implementation checks do not constitute a proof-assistant
formalization of this entire metatheorem.

% Standalone section.  Requires only amsmath and amssymb in the parent document.
% No theorem environments, macros, or bibliography keys are assumed.
% All arithmetic claims in this section are proved here.

\section{From literal-equation induction to full intuitionistic PRA}
\label{sec:eic-completion}\label{app:arithmetic}

This section supplies an unconditional completion theorem. It is adapted
with attribution from the current forward-left draft~\cite{forward}. In the
scalar application used here, every function occurring in the proof is
already a scalar term: predecessor is supplied by
Appendix~\ref{sec:hybrid-bootstrap}, the selector is
$C(n,a,b)=I(n,a;u.b)$, addition is $x+y=I(y,x;u.Su)$, and
truncated subtraction is $d(x,y)=I(y,x;u.Pu)$. All other functions below
are finite compositions of these. Thus this proof does not presuppose
that arbitrary primitive recursion has already been compiled. The QFI
conclusion first applies to the scalar language; the general PR compiler
comes afterward.  Literal-equation
induction, together with the primitive-recursive defining equations,
predecessor coverage, and separation of $0$ and $1$, proves induction for every
quantifier-free formula.  In particular, the proof does not assume a faithful
numerical equality test, stability or decidability of equality, successor
uniqueness, or any commutation law for an iterator.  We construct the equality
test and prove its reflection property for arbitrary elements of the
numerical domain.

\subsection{The exact hypotheses}

Work in ordinary intuitionistic first-order logic with equality and genuine
falsity $\bot$.  The language here is the staged scalar-iterator language of the main text.
All operations in this appendix are the finite scalar terms listed in its
opening paragraph. Only their already proved equations are used; general
primitive recursion is compiled in the following appendix, after QFI.
In particular, predecessor and selector satisfy
\begin{align}
 P(0)&=0, & P(Sx)&=x, \label{eq:eic-predecessor}\\
 C(0,a,b)&=a, & C(Sx,a,b)&=b. \label{eq:eic-selector}
\end{align}
Write $1=S0$.  Besides those scalar equations, use
\begin{align}
 \mathsf{Cov}:&\qquad
       \forall x\,\bigl(x=0\lor x=S(Px)\bigr),
       \label{eq:eic-coverage}\\
 \mathsf{Sep}:&\qquad 0=1\longrightarrow\bot,
       \label{eq:eic-separation}
\end{align}
and induction for each \emph{literal equation} $E(n,\vec p)$:
\begin{equation}
 \forall\vec p\,
 \left[
   \left(E(0,\vec p)\land
       \forall n\bigl(E(n,\vec p)\longrightarrow E(Sn,\vec p)\bigr)
   \right)
   \longrightarrow\forall n\,E(n,\vec p)
 \right].
 \label{eq:eic-induction}
\end{equation}
Here the schematic expression $E$ is a single equality of numerical
terms; the displayed universal quantifiers range only over numbers.
Every description and every instance of the schema is finite external
syntax.

The axiom formulation \eqref{eq:eic-induction} and the usual contextual
induction rule are interchangeable.  For example, if
\[
 \Gamma\vdash E(0,\vec p),
 \qquad
 \Gamma,E(z,\vec p)\vdash E(Sz,\vec p),
\]
where $z$ is absent from $\Gamma$ and the passive parameters, discharge the
second premise's equation and universally introduce $z$.  The axiom then
gives $\Gamma\vdash\forall n E(n,\vec p)$, which can be instantiated at any
target term.  Conversely, assume the antecedent of the axiom and apply the
contextual rule with fresh variables.  This uses the ordinary intuitionistic
logical rules and permits arbitrary, including quantified, formulas in
$\Gamma$.

\subsection{Constructor separation and two elementary tests}

Define the sign description and the countdown description by
\begin{equation}
 \operatorname{sg}(u)=C(u,0,1),
 \qquad
 K(u)=C(u,1,Pu).
 \label{eq:eic-sign-countdown}
\end{equation}
Their defining equations imply
\begin{align}
 \operatorname{sg}(0)&=0,
 &\operatorname{sg}(Sx)&=1,\label{eq:eic-sign-computation}\\
 K(0)&=1,
 &K(Sx)&=x.\label{eq:eic-countdown-computation}
\end{align}

\paragraph{Constructor laws.}
Successor is injective: applying $P$ to $Sx=Sy$ gives $x=y$ by
\eqref{eq:eic-predecessor}.  Successor never equals zero.  Indeed,
$Sx=0$ gives $1=0$ by applying $\operatorname{sg}$ and using
\eqref{eq:eic-sign-computation}; symmetry and
\eqref{eq:eic-separation} give $\bot$.  Consequently $1\ne0$, where
$A\ne B$ abbreviates $A=B\to\bot$.  None of these arguments uses
induction.

\paragraph{The sign test.}
For every $u$,
\begin{align}
 &\operatorname{sg}(u)=0\lor\operatorname{sg}(u)=1,
       \label{eq:eic-sign-range}\\
 &\operatorname{sg}(u)=0\quad\longleftrightarrow\quad u=0,
       \label{eq:eic-sign-zero}\\
 &\operatorname{sg}(u)=1\quad\longleftrightarrow\quad u\ne0.
       \label{eq:eic-sign-one}
\end{align}
To prove the range assertion, split \eqref{eq:eic-coverage}.  The branch
$u=0$ gives sign $0$, and the branch $u=S(Pu)$ gives sign $1$.
The implication $u=0\to\operatorname{sg}(u)=0$ is computation.
For its converse, the successor branch of coverage would make
$\operatorname{sg}(u)=0$ imply $1=0$, a contradiction; in that branch
explosion gives $u=0$ as well.  Similarly, if
$\operatorname{sg}(u)=1$ and $u=0$, computation gives $0=1$.
Conversely, $u\ne0$ excludes the zero branch of coverage, and the
successor branch computes the sign as $1$.

Coverage and \eqref{eq:eic-selector} also give, for all $u,v$,
\begin{equation}
 C(u,v,v)=v.
 \label{eq:eic-selector-idempotence}
\end{equation}
Both coverage branches compute to $v$, so this identity again needs no
induction.

\paragraph{The countdown test and the notation $W$.}
For every $u$,
\begin{equation}
 K(u)=0\quad\longleftrightarrow\quad u=1.
 \label{eq:eic-countdown-zero}
\end{equation}
If $u=1=S0$, \eqref{eq:eic-countdown-computation} gives $K(u)=0$.
Conversely, assume $K(u)=0$ and split coverage.  If $u=0$, then
$K(u)=1$, contradicting separation, and explosion gives $u=1$.
If $u=S(Pu)$, computation gives $K(u)=Pu$; hence $Pu=0$ and
$u=S0=1$.

If a presentation uses a description named $W$ with defining equations
\[
 W(0)=1,\qquad W(Sx)=x,
\]
then it proves $K(u)=W(u)$ for every $u$ already from coverage.
In the zero branch both sides are $1$; in the branch $u=S(Pu)$ both
sides are $Pu$.  Thus $K$ and $W$ are provably equal on the whole domain,
even if they are distinct descriptions.  This equality does not appeal to
standard-model agreement or to uniqueness of an arbitrary recurrence.

\subsection{Four elementary uses of equation induction}

Choose descriptions for addition and truncated subtraction, with recursion
on the second argument:
\begin{align}
 x+0&=x, &x+Sy&=S(x+y),\label{eq:eic-addition}\\
 d(x,0)&=x, &d(x,Sy)&=P(d(x,y)).\label{eq:eic-subtraction}
\end{align}
Using these equations and \eqref{eq:eic-induction}, we now derive
\begin{align}
 0+y&=y,\label{eq:eic-left-zero}\\
 (Sx)+y&=S(x+y),\label{eq:eic-left-successor}\\
 d(Sx,Sy)&=d(x,y),\label{eq:eic-cancel-successors}\\
 d(Sa,a)&=1.\label{eq:eic-diagonal}
\end{align}
Each induction below has a single equation as its induction formula.

For \eqref{eq:eic-left-zero}, induct on $y$.  The base is $0+0=0$.
The step is $0+Sy=S(0+y)=Sy$.

For \eqref{eq:eic-left-successor}, hold $x$ fixed and induct on $y$.
At zero both sides are $Sx$.  At a successor the induction hypothesis
gives
\[
 (Sx)+Sy=S\bigl((Sx)+y\bigr)
       =S\bigl(S(x+y)\bigr)=S(x+Sy).
\]
Combining \eqref{eq:eic-left-successor} with the defining equation for
addition gives the useful rearrangement
\begin{equation}
 (Sy)+r=y+Sr.
 \label{eq:eic-shift-successor}
\end{equation}
Indeed, both sides are $S(y+r)$.

For \eqref{eq:eic-cancel-successors}, hold $x$ fixed and induct on $y$.
The base computes as
\[
 d(Sx,S0)=P(d(Sx,0))=P(Sx)=x=d(x,0).
\]
Applying $P$ to the induction equation gives the step:
\[
 d(Sx,S(Sy))=P(d(Sx,Sy))=P(d(x,y))=d(x,Sy).
\]

Finally, induct on $a$ for \eqref{eq:eic-diagonal}.  The base is
$d(S0,0)=S0=1$.  At a successor,
\[
 d(S(Sa),Sa)=d(Sa,a)=1,
\]
using \eqref{eq:eic-cancel-successors} and the induction hypothesis.

\subsection{An equation invariant that reflects actual equality}

The remaining arithmetic ingredient is the following equation, valid for
all $x,y$:
\begin{equation}
 C\bigl(\operatorname{sg}(d(x,y)),\ x,\ y+d(x,y)\bigr)=x.
 \label{eq:eic-residual-invariant}
\end{equation}
Its meaning is that a nonzero subtraction residual reconstructs the
original argument by addition.  The selector expresses this assertion as
one literal equation, so equation induction applies directly.

Hold $x$ fixed and induct on $y$.  At $y=0$, subtraction gives residual
$x$, and \eqref{eq:eic-left-zero} gives $0+x=x$.  The desired equation
is consequently $C(\operatorname{sg}(x),x,x)=x$, which is
\eqref{eq:eic-selector-idempotence}.

For the step, assume \eqref{eq:eic-residual-invariant} at $y$ and
abbreviate
\[
 r=d(x,y),\qquad r'=Pr=d(x,Sy).
\]
These are local term abbreviations, not additional function symbols.
Split coverage for $r'$.

If $r'=0$, then $\operatorname{sg}(r')=0$, and the successor instance
of \eqref{eq:eic-residual-invariant} reduces immediately to $x=x$.

In the other branch, $r'=S(Pr')$, so constructor separation gives
$r'\ne0$.  Split coverage for $r$.  Its zero branch would imply
$r'=P0=0$, a contradiction, from which the desired step equation follows.
In its successor branch,
\[
 r=S(Pr)=Sr',\qquad \operatorname{sg}(r)=1.
\]
The induction equation therefore computes to $y+r=x$.  Meanwhile,
\eqref{eq:eic-shift-successor} gives
\[
 (Sy)+r'=y+Sr'=y+r=x.
\]
Since $r'$ is a successor, $\operatorname{sg}(r')=1$, so this is
exactly the branch selected in the successor instance of
\eqref{eq:eic-residual-invariant}.  This completes the equation induction.
All case distinctions came from coverage; no induction on an implication
or a disjunction was used.

\paragraph{Reflection at residual $1$.}
We have, for arbitrary $a,b$,
\begin{equation}
 d(Sa,b)=1\quad\longleftrightarrow\quad a=b.
 \label{eq:eic-residual-reflection}
\end{equation}
For the forward direction, substitute $x=Sa$, $y=b$ into
\eqref{eq:eic-residual-invariant}.  If $d(Sa,b)=1$, the sign is $1$
and the selector returns its second branch, giving
\[
 b+1=Sa.
\]
The defining equation for addition gives $b+1=S(b+0)=Sb$.  Thus
$Sb=Sa$, and applying $P$ gives $b=a$, hence $a=b$.
Conversely, $a=b$ and \eqref{eq:eic-diagonal} give
$d(Sa,b)=d(Sa,a)=1$.

Define the equality characteristic by the fixed description
\begin{equation}
 e(a,b)=\operatorname{sg}\bigl(K(d(Sa,b))\bigr).
 \label{eq:eic-equality-characteristic}
\end{equation}
The sign, countdown, and residual reflection lemmas give
\begin{align}
 e(a,b)=0&\quad\longleftrightarrow\quad a=b,
       \label{eq:eic-equality-reflection}\\
 e(a,b)=0&\ \lor\ e(a,b)=1.
       \label{eq:eic-equality-range}
\end{align}
In detail, $e(a,b)=0$ is equivalent to $K(d(Sa,b))=0$ by
\eqref{eq:eic-sign-zero}, which is equivalent to $d(Sa,b)=1$ by
\eqref{eq:eic-countdown-zero}, and then to $a=b$ by
\eqref{eq:eic-residual-reflection}.  The range assertion is
\eqref{eq:eic-sign-range} applied to $K(d(Sa,b))$.

The reflection property concerns the logic's actual equality throughout
the domain.  It was not assumed as part of a numerical comparison symbol,
and its proof did not restrict $a$ or $b$ to standard numerals.
It also yields decidable equality.  In the $e(a,b)=0$ branch,
\eqref{eq:eic-equality-reflection} gives $a=b$.  In the $e(a,b)=1$
branch, an assumption $a=b$ would give $e(a,b)=0$ and therefore $0=1$,
contradicting separation.  Hence
\begin{equation}
 \forall a\,\forall b\,\bigl(a=b\lor(a=b\to\bot)\bigr)
 \label{eq:eic-equality-decidable}
\end{equation}
is a theorem, not an extra hypothesis.

\subsection{Compilation of all quantifier-free formulas}

For each fixed quantifier-free formula $A(\vec p)$, construct a term
$b_A(\vec p)$ by external recursion on its finite syntax.  Take equality,
$\bot$, $\land$, $\lor$, and $\to$ as the primitive formula constructors;
negation is implication into $\bot$.  Set
\begin{align}
 b_{s=t}&=e(s,t), & b_{\bot}&=1,\label{eq:eic-boolean-atoms}\\
 b_{A\land D}&=C(b_A,b_D,1),\label{eq:eic-boolean-and}\\
 b_{A\lor D}&=C(b_A,0,b_D),\label{eq:eic-boolean-or}\\
 b_{A\to D}&=C(b_A,b_D,0).\label{eq:eic-boolean-implies}
\end{align}
Every term here is a finite composition of fixed primitive-recursive
descriptions; if desired, projections and composition give it a single
description symbol.  No internal quantification over formulas or programs
is involved.

An external induction on the syntax of $A$ proves, within the arithmetic
theory,
\begin{equation}
 \bigl(b_A=0\lor b_A=1\bigr)
 \quad\text{and}\quad
 \bigl(b_A=0\longleftrightarrow A\bigr).
 \label{eq:eic-boolean-correctness}
\end{equation}
Here are the full connective checks.  In every case $b_A=1$ implies
$\neg A$: if $A$ held, the already established implication
$A\to b_A=0$ would contradict separation.

\paragraph{Atomic formulas and falsity.}
For an equation, both claims are
\eqref{eq:eic-equality-reflection}--\eqref{eq:eic-equality-range}.
For falsity the range assertion is $1=0\lor1=1$, with its right disjunct
proved by reflexivity.  Also $1=0\to\bot$ follows from separation and
symmetry, while $\bot\to1=0$ follows by explosion.  Thus $1=0$ is
equivalent to $\bot$, as required.

\paragraph{Conjunction.}
Split $b_A=0\lor b_A=1$.  In the zero branch, $A$ holds and
$b_{A\land D}=b_D$.  Its range follows from the range of $b_D$,
and its being zero is equivalent to $D$, hence to $A\land D$ under
the known $A$.  In the one branch, $\neg A$ holds and
$b_{A\land D}=1$.  Its range is immediate.  If it were zero,
separation would give $\bot$, hence $A\land D$; conversely an
assumption $A\land D$ contradicts $\neg A$, and explosion gives
$b_{A\land D}=0$.

\paragraph{Disjunction.}
In the branch $b_A=0$, the formula $A$ holds and
$b_{A\lor D}=0$.  The range is immediate, and both $b_{A\lor D}=0$
and $A\lor D$ hold, proving their equivalence.  In the branch
$b_A=1$, we have $\neg A$ and $b_{A\lor D}=b_D$.  Its range follows
from the range of $b_D$.  If its value is zero, then $D$ and hence
$A\lor D$.  Conversely, from $A\lor D$, the $A$ alternative
contradicts $\neg A$ and yields $b_D=0$ by explosion, while the
$D$ alternative yields $b_D=0$ directly.  This proves the required
equivalence in both alternatives.

\paragraph{Implication.}
In the branch $b_A=0$, the formula $A$ holds and
$b_{A\to D}=b_D$.  Its range follows from the range of $b_D$.
If its value is zero, then $D$, which gives $A\to D$.  Conversely,
$A\to D$ and the known $A$ give $D$, hence $b_D=0$.
In the branch $b_A=1$, we have $\neg A$ and
$b_{A\to D}=0$.  Its range is immediate.  To prove $A\to D$,
assume $A$; this contradicts $\neg A$, and explosion gives $D$.
Thus both $b_{A\to D}=0$ and $A\to D$ hold in this branch, and
their equivalence follows.  This argument uses the proved two-value
range of the finite characteristic.  It does not use excluded middle
for arbitrary formulas or change the surrounding intuitionistic logic.

This completes the external syntactic induction proving
\eqref{eq:eic-boolean-correctness} for every quantifier-free formula.

\subsection{Recovery of induction and the full theory}

Let $A(n,\vec p)$ be any fixed quantifier-free formula.  Suppose an
arbitrary retained context $\Gamma$ proves its base and its induction
step:
\[
 \Gamma\vdash A(0,\vec p),
 \qquad
 \Gamma,A(z,\vec p)\vdash A(Sz,\vec p),
\]
with the usual freshness condition on $z$.  By
\eqref{eq:eic-boolean-correctness}, the base gives
$\Gamma\vdash b_A(0,\vec p)=0$.  Under the additional assumption
$b_A(z,\vec p)=0$, the same equivalence gives $A(z,\vec p)$,
the given step gives $A(Sz,\vec p)$, and the equivalence gives
$b_A(Sz,\vec p)=0$.  Apply literal-equation induction to
$b_A(n,\vec p)=0$ and then apply its reflection implication.
The result is $\Gamma\vdash A(t,\vec p)$ for every target term $t$.
Thus the theory proves the full contextual quantifier-free induction
rule, including in arbitrary quantified contexts.

\paragraph{Completion theorem.}
In the common language of fixed primitive-recursive descriptions, the
theory
\[
 \mathsf{PRE}+\mathsf{Cov}+\mathsf{Sep}+\mathsf{EI}
\]
has exactly the same sequents as ordinary first-order intuitionistic
primitive recursive arithmetic with all quantifier-free induction
instances.  The forward simulation follows from the recovered
quantifier-free induction, successor injectivity, and successor separation;
all primitive-recursive defining equations and all logical rules are
shared.  For the reverse simulation, literal-equation induction is an
instance of quantifier-free induction, and successor separation gives
$\mathsf{Sep}$.  Quantifier-free induction proves $\mathsf{Cov}$:
at zero choose its left disjunct, and at a successor use
$Sx=S(P(Sx))$, which follows from $P(Sx)=x$, for its right disjunct.

Both interpretations therefore use the full domain, defined by $x=x$,
and preserve every function symbol, equality, logical connective, and
quantifier.  Both composites have the identity isomorphism, with graph
$x=y$.  The theories are consequently identity-domain bi-interpretable
and equiconsistent.  This completion theorem is unconditional once the
literal-equation induction rule has been obtained; it has no dependency
on a classical successor-uniqueness reduction.

\section{Pairing, primitive recursion, and internal numerical composites}
\label{sec:iterator-coding}
This appendix is adapted with attribution from the progress-closure
arithmetic development~\cite{progress}. Its prerequisites have now been
proved: scalar predecessor, arithmetic reflection, and full quantifier-free
induction. Pairing is a derived proof tool; it is not part of any checker.

\subsection{Pairing from a bounded square-root search}

Once QFI is available, a familiar square-based injection gives a
particularly direct compiler. Write $a\le b$ for $d(a,b)=0$.
The usual linear-order and semiring identities follow by fixed
quantifier-free inductions from the successor, addition, and monus
equations. Multiplication itself is the scalar term
\[
                        a\cdot b=I(b,0;u.u+a).
\]
Only primitive-recursive arithmetic is involved in the following
bounded root search; there is no minimization over an unbounded domain.

For a fixed parameter $n$, put
\[
 U_n(v)=C(d((Sv)^2,n),Sv,v),\qquad
 R_n(i)=I(i,0;v.U_n(v)),\qquad q(n)=R_n(n).
\]
The update increments $v$ exactly when $(v+1)^2\le n$, and
otherwise retains it. QFI in $i$, with $n$ fixed, proves
\[
 R_n(i)\le i,\qquad R_n(i)^2\le n,\qquad
 R_n(i)=i\ \vee\ n<(S(R_n(i)))^2.              \tag{Root invariant}
\]
The base has $R_n(0)=0$. If the update increments, its test
contradicts the strict-upper-bound alternative in the third invariant.
Thus the old state equals $i$, and the new state equals $Si$;
its square bound is the increment test. If the update does not increment, the
strict upper square bound holds and the state remains fixed.
An already established strict upper square bound forces every
subsequent update to remain fixed. These are exactly the cases
needed for the three displayed invariants.

At $i=n$, either the strict upper square bound already holds or
$R_n(n)=n$, in which case $n<(n+1)^2$ supplies it. Thus
\[
                        q(n)^2\le n<(q(n)+1)^2.           \tag{Root}
\]
Monotonicity of squares shows uniqueness of a number satisfying
Root. All formulas just used are quantifier free in the scalar
language, with fixed $n$ during the induction.

Now define
\[
 \pi(a,b)=(a+b)^2+a,\qquad
 L(n)=d(n,q(n)^2),\qquad H(n)=d(q(n),L(n)).
\]
For $k=a+b$, one has $k^2\le k^2+a<(k+1)^2$ since $a\le k$.
Root uniqueness gives $q(\pi(a,b))=a+b$, and therefore
\[
                    L(\pi(a,b))=a,\qquad H(\pi(a,b))=b.
\]
Surjectivity of $\pi$ is unnecessary. The original numerical domain
is never replaced by the range of this coding; pairing is used only
to store the finite state of a fixed primitive recursion. This
standard square pairing is a convenient reconstruction here, not
a claimed new arithmetic basis.

\subsection{Compiling each primitive recursion}

Suppose the fixed target description is
\[
 f(0,\vec p)=g(\vec p),\qquad
 f(Sn,\vec p)=h(n,f(n,\vec p),\vec p).
\]
Assume by external structural induction that the earlier $g,h$
already have scalar translations $g^*,h^*$. Define
\[
 W(u,\vec p)=\pi(SL(u),h^*(L(u),H(u),\vec p)),
\]
\[
 V(n,\vec p)=I(n,\pi(0,g^*(\vec p));u.W(u,\vec p)),
 \qquad f^*(n,\vec p)=H(V(n,\vec p)).
\]
The update has only its state variable bound, with every $p_i$
fixed. Equation induction proves $L(V(n,\vec p))=n$.
The two projection equations and scalar computation then prove
both defining equations for $f^*$ internally. Initial functions
and composition translate directly. This completes the translation
of every fixed finite PR description without adding a universal
evaluator.

\subsection{Why both numerical round trips are internal}

Two fixed scalar iterators with the same seed and provably equal
updates are internally equal at all counters: use literal-equation
induction comparing their values. Their base agrees, and their
successor equations plus congruence prove the step. The same
uniqueness argument works for ordinary primitive recursion.

For an ordinary PR description, induct externally on its finite
construction. The double translation satisfies its original
defining equations, so internal equation induction identifies it
with the original operation. For a native scalar description,
earlier translated updates agree by the external induction
hypothesis; scalar uniqueness identifies the double-translated
iterator. These proofs hold at arbitrary internal counters.

Thus the domain relation $x=y$ is provably a total bijection
preserving every fixed numerical operation in both composites.
The formula constructions and round trips in
Section~\ref{sec:choice-introduction} give the corresponding formula
composites. Both interpretations preserve falsity, taking native
$0=1$ to an ordinary contradiction and ordinary $\bot$ to native
$0=1$. Finite proof translations therefore establish equiconsistency.

\section{Arithmetic strength and both ordinary interpretation composites}
\label{sec:ordinary-bridge}

We spell out the arithmetic bridge to avoid identifying standard numerical
agreement with bi-interpretability. First elaborate every fixed iterator
into an ordinary function symbol. At stage zero the symbols are $0,S$.
Given a fixed earlier-stage term $h(u,\vec p)$, introduce
$i_h(n,a,\vec p)$ with the two equations
\[
i_h(0,a,\vec p)=a,\qquad
i_h(Sn,a,\vec p)=h(i_h(n,a,\vec p),\vec p).
\]
Only finite descriptions occur; the domain contains no code for a
universal evaluator. Formal parameters are distinct when a description
is declared, while actual arguments can coincide.

Let $T_I$ be ordinary intuitionistic first-order logic with actual
equality and genuine falsity, these staged equations, contextual
induction for literal equations, and
\[
\forall x\,(x=0\vee\exists y\,x=Sy),
\qquad (0=S0)\to\bot.
\]
The second sentence is the falsity bridge. It must be stated: positive
numerical equations by themselves admit a singleton structure. Native
explosion from $0=S0$ and a translation of ordinary falsity to that
equation do not authorize silently dropping the bridge from $T_I$.

The native-to-ordinary translation is
\[
(s=t)^\circ=(s=t),\qquad
(\mathsf D_{x,y}B)^\circ=
\exists x\bigl(B^\circ(x,0)\to\forall y B^\circ(x,y)\bigr).
\]
A clause $\langle A;y.B\rangle$ becomes $A^\circ\to\forall yB^\circ$.
For the auxiliary closure used by the compact checkers,
$(\mathsf Q_yB)^\circ=B^\circ(0)\to\forall yB^\circ$.
The interpretation domain is all elements, defined by $x=x$;
equality is actual equality.

Case DI is sound because numerical coverage combines its retained base
with its successor proof. Progress DI is sound by literal-equation
induction under a temporarily assumed base. Use, DE and transport are
ordinary logical inferences under the translations. Conversely the
derived logical operations, equation induction, cases and explosion in
Section~\ref{sec:choice-introduction} recover the rules and axioms of
$T_I$. An auxiliary closure is interderivable with
$\mathsf D_{v,y}B(y)$ for unused $v$: in one direction use case DI with
the closure and its base; in the other eliminate the paired formula
and return its supplied closure. Thus formulas and context clauses have
the stated meanings in both presentations.

\subsection{The arithmetic dependency order}

The scalar update has no recursion-counter argument. We do not conceal
one in an encoding. The inherited internal proofs, reproduced in the
appendices, establish the required arithmetic in this order:
\begin{enumerate}
\item Using equation induction and numerical cases, construct a fixed
scalar predecessor term and prove $P0=0$ and $P(Sn)=n$.
Existential coverage then gives $x=0\vee x=S(Px)$: in the successor
case $x=Sy$, apply $P$ to obtain $Px=y$ and substitute.
\item Construct truncated subtraction and a numerical characteristic
of actual equality. Prove its reflection on the entire original domain.
\item Compile each fixed quantifier-free formula into its characteristic
equation and recover every contextual quantifier-free induction instance.
\item With that induction available, construct pairing and projections
by a bounded square-root search and compile every fixed general
primitive-recursive description into a scalar iterator term.
\end{enumerate}
No pairing or faithful equality detector is assumed to obtain
predecessor. No appeal to standard-model correctness replaces an
internal equality. Every induction fixes its passive parameters;
diagonal substitution occurs only into a previously proved identity.

For a target recursion
\[
f(0,\vec p)=g(\vec p),\qquad
f(Sn,\vec p)=h(n,f(n,\vec p),\vec p),
\]
let $\pi,L,H$ be the proved pairing and projections. The compiler uses
\[
V(n,\vec p)=I\bigl(n,\pi(0,g^*(\vec p));
u.\pi(SL(u),h^*(L(u),H(u),\vec p))\bigr),
\qquad f^*(n,\vec p)=H(V(n,\vec p)).
\]
Equation induction proves $L(V(n,\vec p))=n$; the two projection
identities then prove the target defining equations. The construction
is external recursion on each finite PR description, not an internal
quantification over arbitrary functions.

\begin{theorem}[Whole-domain bi-interpretability and equiconsistency]
The ordinary realization $T_I$ and full first-order iPRA are
bi-interpretable. In both theories $x=y$ defines the isomorphism from
the original domain to its double interpretation. The native Choice-DI
calculus is therefore a faithful proof presentation of a theory
bi-interpretable and equiconsistent with iPRA.
\end{theorem}
\begin{proof}
Interpret an iterator in iPRA by primitive recursion that ignores its
counter argument in the update. iPRA proves its equations, literal
equation induction, coverage and separation. In the reverse direction
use the scalar compiler. The predecessor and equality-reflection proofs
supply the successor axioms and every quantifier-free induction instance.
All ordinary logical rules are preserved.

For two iterators with the same seed and provably equal updates,
equation induction on the equality of their values proves uniqueness at
every counter, with all parameters fixed. The corresponding argument
applies to ordinary primitive recursion. Induct externally on a finite
description: its double translation satisfies its original defining
equations, so internal uniqueness identifies the two functions. This
also proves coherence when substitution changes the staged syntactic
description of an update.

The domain and equality are unchanged; hence the identity graph is
internally total, bijective and preserves every interpreted operation.
The displayed ordinary meanings of the logical macros and the paired
binder give both formula composites by external induction on finite
formula syntax. Clauses are similarly equivalent to their conditional
universal meanings. The explicit falsity bridge preserves contradiction
in both directions. Effective translations of finite proofs give
equiconsistency, without a claim that either theory proves its own
consistency.
\end{proof}


\section{Attribution of the raw-input audit}
The following raw-input analysis is reproduced from the neutral-major
draft at the pinned snapshot.\cite{neutral} Its occurrences of ``new''
refer to that earlier protocol. Section~\ref{sec:kernel} proves the
additional lexical-scope and branch-folding changes used here. The
minimal-successful-trace argument below quantifies over arbitrary starting
floors, so it continues to apply when major proofs allocate only local
names. The original unmodified source is preserved in \file{reference/}.
% Independent audit fragment for the corrected array-truth interface.
% Reviewed files: kernel_notes/array_truth_checker{,_bytes,_joint}.py.
% IMPORTANT: the final mode comparison is with u, not s.
% Requires only ordinary LaTeX plus amsmath/amssymb.

\section{Why the absent major can be a proved empty array}
\label{sec:neutral-major-audit}

This section audits the new control-flow change on raw input, including
noncanonical JSON.  It uses the inherited sorted arity decoder and
semantic-opening lemmas for the numerical and logical rules.  It does not
replace those lemmas with finite testing or claim proof-assistant
certification.

\subsection{The exact change}

The initial context is the one-element list \texttt{[[]]}.  Its sole
formula, the empty array, denotes $0=0$ under the inherited decoder.
At every non-hypothesis command the checker now executes
\begin{verbatim}
    m, s, u, t = q()
    c(m, g)
\end{verbatim}
without testing whether \texttt{m} is nonempty.  A canonical command with
no substantive major puts \texttt{[]} in this field.  Its premise closes
immediately by hypothesis, without consuming a command or changing the
fresh-name floor.

The final choice of the retained assumption uses
\begin{verbatim}
    (m or len(a) % 2 == u) and k
\end{verbatim}
instead of comparing the arity parity with \texttt{m}.  In an
introduction command, \texttt{u=0} selects equation progress and
\texttt{u=[]} selects the zero-case premise.  This reuse is legitimate
because introduction is reached only when both \texttt{m} and
\texttt{u} are false-valued.  Computation still requires a truthy
initializer \texttt{u}; its missing literal-zero cases have the
inherited certified-zero repair.

The source is deliberately left with the old matcher
\begin{verbatim}
    m in ([s, u], not s and [t] * 3)
\end{verbatim}
and no added \texttt{m and} guard.  As an isolated Python expression this
would be unsafe: if \texttt{s} is truthy, its second alternative is
\texttt{False}, which compares equal to numeric zero.  The following
whole-program invariant proves that this dangerous major cannot reach
the matcher.

\subsection{Two facts about the root opener}

Call a Python value \emph{numeric} here when it is an integer, a float,
or a Boolean returned by the JSON loader.  The finite Python execution
also admits the loader's optional nonfinite float values.  No custom
numeric class is supplied by a certificate.

Write $O_k(r;t)$ for the opener on the numeric value $r$, with image $t$
and its default formula-root depth zero.  Direct evaluation of the
program gives
\[
 O_k(r;t)=
 \begin{cases}
 r,&k\le r<0,\\
 0,&\text{otherwise}.
 \end{cases}
 \tag{Root scalar}
\]
In particular this result is numeric and independent of the image $t$.
The formula covers negative allocated names and negative fractional
names without identifying different numerical values.  Signed zero and
Booleans have Python's ordinary numerical equality behavior.  A NaN
fails the interval test and returns zero; the relevant infinite values
also return zero.  None of these results is a list.

A second fact is that a successful root opening of a list returns a
list.  Under either strict list predicate, no other input type can
return a numeric result except a numeric input.  Under the generic
iterable opener, a successful iterable traversal also returns a list;
if traversal fails, the numerical comparison cannot succeed on a list,
string, dictionary, or null.  Thus in every profile:
\begin{quote}
If a root opening returns a numeric result, its argument was numeric;
if it succeeds on a list, its result is a list.
\end{quote}
These statements concern finite successful evaluations.  They assert
neither arbitrary revalidation idempotence at ignored annotations nor
unlimited-stack termination for the generic iterable profile.

\subsection{No numeric goal can be proved from list-valued assumptions}

\paragraph{Lemma.}
For any initial integer fresh-name floor and any finite remaining
certificate, there is no finite successful invocation
\texttt{c(a,g)} whose raw goal \texttt{a} is numeric and whose every
context element is a list.

\emph{Proof.}
Suppose otherwise, and select such a successful invocation with the
fewest calls to \texttt{c} in its completed recursive trace.  Its first
hypothesis check cannot close: the root opening of its numeric goal is
numeric, and no numerical JSON value equals a Python list.

The invocation must therefore consume a command and successfully return
from \texttt{c(m,g)}.  If \texttt{m} were numeric, this would be a
strictly shorter successful invocation of the forbidden kind.  Thus the
major is not numeric.

In either strict profile, a successfully checked nonnumeric major is a
list.  In the generic profile an empty string or an empty dictionary can
also open to the empty list.  A nonempty dictionary can occasionally
have its iterable alias proved, but it cannot occur as a truthy major
in the selected completed parent invocation: the subsequent expression
\texttt{m * 3} raises a type error.  A nonempty string has no finite
successful opening.  Consequently every truthy major that survives to
the continuation exports a list, and every surviving false-valued
major is an empty list or one of those harmless empty iterable aliases.

Consider the possible change of goal.  If the matcher does not fire,
the goal stays numeric.  If it fires, its target test requires
\texttt{i(t,u) == a}.  The second root-opening fact therefore implies
that \texttt{t} is numeric.  By (Root scalar), \texttt{i(t,s)} and
\texttt{i(t,u)} are the same root result.  Both the Boolean-expression
and the explicit-conditional implementations consequently retain a
numeric goal.  The Boolean implementation may retain the original raw
zero alias when the transported result is false-valued; that makes no
difference to this assertion.

If the major is truthy, the generated hypothesis is the list
\texttt{m * 3}.  If the major is false-valued and \texttt{u} is truthy,
the computation expression constructs a two-element list.  A successful
opening of either hypothesis is a list.  The final recursive call then
has a numeric root-opened goal and an entirely list-valued context.  It
is a proper subtrace of the selected invocation, contradicting its
minimality.

The only remaining case is introduction.  Its expression
\texttt{i(*a,t,2)} must unpack the numeric raw goal as an iterable,
which raises a type error.  That case cannot belong to a successful
trace.  All possibilities have been excluded.\hfill$\square$

The minimization ranges over all starting floors and stacks.  It is
therefore unaffected by a recursively checked major decreasing the
global floor before its parent resumes.  The argument uses finite
completed traces, not a presumption that recursive checking always
terminates.

\subsection{Every reached command has a list-valued context}

Initially the only assumption is the list \texttt{[]}.  Suppose a
command is reached with list-valued assumptions and its major check
returns.  The lemma excludes a numeric major.  A truthy surviving major
exports a list; a computation constructs a list.  In introduction the
assignment is
\begin{verbatim}
    h = a = i(*a, t, 2)
\end{verbatim}
in mathematical order.  If its result is numeric, the later
\texttt{len(a)} raises a type error before the continuation context is
constructed, since this branch has false-valued \texttt{m}.  Such a
scalar matrix does \emph{not} insert a scalar-zero assumption.  If the
assignment yields a list, successful opening again yields a list.

Both recursive premises thus receive list-valued contexts whenever
execution reaches them.  The expression \texttt{g + [new]} constructs a
new list; a premise cannot mutate its caller's context.  Induction on
the reached calls proves the invariant.  The tuple-based explanatory
variant has the corresponding immutability property.

It follows that no numeric zero, Boolean \texttt{False}, floating zero,
truthy numeric scalar, or newly allocated numeric name can be a
successfully verified major at a reached command.  In particular a
closed scalar-zero goal is rejected by this interface.  The canonical
empty-array truth is accepted instead; mathematical completeness is
claimed for canonical encodings, not for every redundant truth alias.

\subsection{The unguarded matcher is safe on every reached continuation}

If \texttt{s} is false-valued, both matcher alternatives are nonempty
lists, so no false-valued major can match.  If \texttt{s} is truthy, the
second alternative is \texttt{False}.  A false-valued major reaching
this point is an empty list, or in the generic profile an empty string
or dictionary.  None equals \texttt{False}; a numeric major has already
been excluded.  Therefore, at every point reached after a successful
major check, the matcher has exactly the same truth value as
\begin{verbatim}
    m and m in ([s, u], not s and [t] * 3)
\end{verbatim}
even though this identity is false on unrestricted Python tuples of
values.  No sentinel and no extra source guard are needed.

The proof depends on the initial context being \texttt{[[]]}, not
\texttt{[0]}, and on premise contexts being freshly constructed.
Silently changing either detail would invalidate this argument.

\subsection{Why the relocated induction mode must be the initializer field}

The corrected selector compares arity parity with \texttt{u}.  In an
introduction branch this field is necessarily false-valued.  No
false-valued JSON value compares equal to one.  Thus a meaningful
compound matrix, of odd arity one or three, cannot receive a progress
assumption.  A literal equation has even arity two and receives
progress precisely for a zero-like mode such as \texttt{0}; an empty
array mode chooses the zero instance.  Surplus arities retain their
inherited constant interpretations.  Scalar matrices cannot survive
the shape check discussed above.

Comparing instead with the unused \texttt{s} field would be wrong.
That field is unrestricted in introduction, so choosing
\texttt{s=1} would enable progress for odd-arity compound matrices.
The distinction is semantic, not merely a preferred serialization.

In the computation branch \texttt{u} is truthy.  Its new comparison can
choose a fresh-name opening for some noncanonical scalar initializer
aliases.  This is handled by the inherited zero-instance discharge
lemma, not by claiming that every generated open hypothesis is
uniformly true.  If the generated hypothesis is $H(\vec b)$, its zero
completion $H(\vec0)$ is a computation theorem, while the old-floor
validation makes the finite family of newly admitted names $\vec b$
fresh for the retained conclusion and context.  Substituting zero for
those names and cutting in the computation theorem discharges the
extra hypothesis.  Negative fractional names in the newly admitted
interval are included separately in this finite substitution.

\subsection{Conservativity and canonical completeness}

The initial empty-array formula is already a theorem.  For example the
permitted nonempty-seed computation gives
\[
 C=0,\qquad C:=I(S(S0),S0;v.0).
\]
Use this equation as both the equality major and the premise of
transport through $F(w)\equiv(w=0)$ to obtain $0=0$.  Adding that fixed
theorem to a proof context is conservative by substitution of its
finite derivation for hypothesis uses.  No new object-level truth
predicate or logical axiom has been introduced.

For canonical certificates, substantive major commands retain their
old fields.  A computation uses the absent-major field \texttt{[]} and
its existing nonempty initializer.  An introduction uses absent major
\texttt{[]}, with \texttt{u=0} for equation progress and
\texttt{u=[]} for cases; its witness remains in \texttt{t}.
The added empty major closes without a command or name allocation, so
the remaining canonical proof simulation is unchanged.  First applying
the existing certified-zero repair supplies every unrestricted
computation instance.  This gives a finite certificate for every
canonical derivation of the inherited native calculus.

Conversely, the invariant makes each reached matcher a legitimate
guarded matcher.  The remaining commands use the inherited
transport, closure, paired-introduction/elimination, and zero-instance
computation arguments, with the corrected equation-only mode test.
Discharging the initial $0=0$ theorem yields a derivation in that same
calculus.  Its established whole-domain iPRA interpretations therefore
transfer without changing the numerical domain, equality, or object
language.  The new result concerns a proof-checker presentation, not a
new primitive-recursive arithmetic foundation.

\subsection{Raw JSON types and the generic profile}

The strict profiles recognize only Python lists as containers.  The
short predicate \texttt{a * 0 == []} has this effect for JSON values:
lists pass; strings and numbers give different zero products;
dictionaries and null raise.  Its numerical fallback rejects strings.
The root invariant covers all numbers and Booleans that reach it,
including nonfinite values accepted by the default Python loader.

The generic profile has the inherited additional aliases.  An empty
string and an empty dictionary open to \texttt{[]}; a dictionary with
only the empty-string key opens to \texttt{[[]]}, with its value ignored.
A truthy dictionary cannot survive export as a major because
\texttt{m * 3} raises.  Nonempty strings and nonempty dictionary keys
have no finite successful opening.  Actual CPython rejects them through
its recursion limit and subsequent type failure, while an ideal
unlimited-stack execution can diverge.  The raw soundness argument
above covers every finite successful run of this profile; it does not
assert ideal termination on every malformed finite JSON input.

Larger arrays remain arrays throughout opening, irrespective of their
ignored annotations or depth increments.  They therefore do not
threaten the scalar-versus-list invariant.  Their formula meanings and
freshness behavior remain subject to the inherited sorted-decoder
proof.  Physical time, memory, and stack limits can reject sufficiently
large genuine proofs in every profile; the completeness claim has the
same adequate-resource, canonical-proof interpretation as its
predecessors.


\begin{thebibliography}{9}
\bibitem{source} Vladimir Reshetnikov, \emph{ProveIt}, frozen draft corpus
at commit \file{b9c4d74a0018bdafc5bc06e755a0145f4ed2b8d6}.
\url{https://github.com/VladimirReshetnikov/ProveIt/tree/b9c4d74a0018bdafc5bc06e755a0145f4ed2b8d6/docs/drafts}.
Retrieved and byte-verified 10 November 2026.
\bibitem{maintained} Vladimir Reshetnikov, \emph{Concise presentations of
intuitionistic primitive recursive arithmetic}, maintained project README,
commit \texttt{89f5bdfe}, 10 November 2026.
\url{https://github.com/VladimirReshetnikov/ProveIt/blob/89f5bdfef815cd2b2c5422978896edc9c819af3f/Logic/PrimitiveRecursiveArithmetic/ConciseIPRA/README.md}.
\bibitem{neutral} \emph{Equality from monus and neutral major premises},
\file{ipra_monus_and_neutral_major_2026-11-10.zip}, in the pinned corpus.
Source of the neutral-major protocol and its raw-input audit.
\bibitem{selfseed} \emph{Scalar self-seeded arithmetic and the iPRA checker
frontier}, \file{ipra_selfseed_scalar_2026-11-10.zip}, in the pinned corpus.
Source of the complete self-seeded presentation, lexical profiles,
and the collected inherited arithmetic appendices.
\bibitem{arity} \emph{Arity calculus for iPRA},
\file{ipra_arity_calculus_2026-11-10.zip}, in the pinned corpus.
\bibitem{certified} \emph{Certified zero},
\file{ipra_certified_zero_2026-11-10.zip}, in the pinned corpus.
\bibitem{erasure} Inherited scalar-predecessor development,
\file{reference/scalar_predecessor.tex} in the self-seeded archive
of reference~\cite{selfseed}. The exact source hash is recorded in
\file{evidence/mathematical_provenance.json}.
\bibitem{forward} Inherited equality-reflection development,
\file{reference/equality_reflection.tex} in reference~\cite{selfseed}.
The presentation here explicitly restricts its function inventory to the
already constructed scalar terms before general PR compilation.
\bibitem{progress} Inherited pairing and primitive-recursion development,
\file{reference/primitive_recursion.tex} in reference~\cite{selfseed}.
\end{thebibliography}
\end{document}
